%%%%%%%%%%%%%%%%%%%%%%%%%%%%%%%%%%%%%%%%%%%%%%%%%%%%%%%%%%%%
% 8.15 TIME SERIES ANALYSIS
% The Composition of Advanced Mathematics
%%%%%%%%%%%%%%%%%%%%%%%%%%%%%%%%%%%%%%%%%%%%%%%%%%%%%%%%%%%%

\section{Time Series Analysis}
\label{sec:time-series-analysis}

\prerequisite{
Probability, random variables, expectation and variance, covariance and
correlation, regression, statistical estimation, hypothesis testing,
stochastic processes, linear algebra, and basic differential and
difference equations.
}

\learningobjectives{
By the end of this section, the reader should be able to:
\begin{itemize}
    \item explain what distinguishes time series data from ordinary
          independent observations;
    \item identify trend, seasonality, cycles, and irregular variation;
    \item define weak stationarity;
    \item define autocovariance and autocorrelation functions;
    \item interpret lag;
    \item compute sample autocovariances and autocorrelations;
    \item understand white-noise processes;
    \item understand random walks;
    \item distinguish stationary and nonstationary time series;
    \item understand differencing;
    \item define autoregressive models;
    \item define moving-average models;
    \item define ARMA and ARIMA models;
    \item understand stationarity conditions for autoregressive models;
    \item understand invertibility for moving-average models;
    \item derive basic mean and variance properties of AR models;
    \item understand the Yule--Walker equations;
    \item interpret ACF and PACF patterns;
    \item understand seasonal ARIMA models conceptually;
    \item understand exponential smoothing;
    \item understand Holt and Holt--Winters methods conceptually;
    \item understand forecasting and forecast uncertainty;
    \item distinguish one-step and multi-step forecasts;
    \item understand residual diagnostics;
    \item recognize unit-root behavior;
    \item understand deterministic and stochastic trends;
    \item understand transformations such as logarithms;
    \item understand intervention and transfer-function ideas
          conceptually;
    \item recognize state-space models;
    \item understand the Kalman-filter idea conceptually;
    \item distinguish univariate and multivariate time series;
    \item recognize vector autoregression;
    \item understand cointegration conceptually;
    \item recognize frequency-domain analysis;
    \item understand spectral density conceptually;
    \item connect time series with stochastic processes, forecasting,
          regression, and dynamical systems.
\end{itemize}
}

%%%%%%%%%%%%%%%%%%%%%%%%%%%%%%%%%%%%%%%%%%%%%%%%%%%%%%%%%%%%
\subsection{What Is a Time Series?}
\label{subsec:what-is-time-series}
%%%%%%%%%%%%%%%%%%%%%%%%%%%%%%%%%%%%%%%%%%%%%%%%%%%%%%%%%%%%

A time series is a sequence of observations indexed by time.

A discrete-time series may be written

\[
\boxed{
X_1,
X_2,
\ldots,
X_T.
}
\]

Examples include:

\[
\boxed{
\text{daily temperature},
}
\]

\[
\boxed{
\text{monthly sales},
}
\]

\[
\boxed{
\text{hourly sensor measurements},
}
\]

or

\[
\boxed{
\text{annual economic indicators}.
}
\]

%%%%%%%%%%%%%%%%%%%%%%%%%%%%%%%%%%%%%%%%%%%%%%%%%%%%%%%%%%%%
\subsection{Why Time Order Matters}
\label{subsec:why-time-order-matters}
%%%%%%%%%%%%%%%%%%%%%%%%%%%%%%%%%%%%%%%%%%%%%%%%%%%%%%%%%%%%

In ordinary i.i.d. sampling, rearranging observations does not change
the probabilistic structure.

In a time series, order is essential.

For example,

\[
X_t
\]

may depend strongly on

\[
X_{t-1}.
\]

Thus

\[
\boxed{
\text{time dependence}
}
\]

is a defining feature of time series analysis.

%%%%%%%%%%%%%%%%%%%%%%%%%%%%%%%%%%%%%%%%%%%%%%%%%%%%%%%%%%%%
\subsection{Time Series as a Stochastic Process}
\label{subsec:time-series-stochastic-process}
%%%%%%%%%%%%%%%%%%%%%%%%%%%%%%%%%%%%%%%%%%%%%%%%%%%%%%%%%%%%

A time series can be viewed as one observed realization of a stochastic
process

\[
\boxed{
\{X_t:t\in\mathbb{Z}\}.
}
\]

The index \(t\) represents discrete time.

The probabilistic model describes the joint behavior of

\[
X_t,
X_{t-1},
X_{t-2},
\ldots
\]

across time.

%%%%%%%%%%%%%%%%%%%%%%%%%%%%%%%%%%%%%%%%%%%%%%%%%%%%%%%%%%%%
\subsection{Components of a Time Series}
\label{subsec:time-series-components}
%%%%%%%%%%%%%%%%%%%%%%%%%%%%%%%%%%%%%%%%%%%%%%%%%%%%%%%%%%%%

A time series may contain several structural components:

\[
\boxed{
\text{trend},
}
\]

\[
\boxed{
\text{seasonality},
}
\]

\[
\boxed{
\text{cycles},
}
\]

and

\[
\boxed{
\text{irregular variation}.
}
\]

These components may overlap.

%%%%%%%%%%%%%%%%%%%%%%%%%%%%%%%%%%%%%%%%%%%%%%%%%%%%%%%%%%%%
\subsection{Trend}
\label{subsec:time-series-trend}
%%%%%%%%%%%%%%%%%%%%%%%%%%%%%%%%%%%%%%%%%%%%%%%%%%%%%%%%%%%%

A trend is a long-term change in the level of a series.

Examples include:

\[
\boxed{
\text{steady growth},
}
\]

\[
\boxed{
\text{steady decline},
}
\]

or

\[
\boxed{
\text{long-run nonlinear change}.
}
\]

A simple deterministic linear trend is

\[
\boxed{
X_t
=
\beta_0+\beta_1t+\varepsilon_t.
}
\]

%%%%%%%%%%%%%%%%%%%%%%%%%%%%%%%%%%%%%%%%%%%%%%%%%%%%%%%%%%%%
\subsection{Seasonality}
\label{subsec:time-series-seasonality}
%%%%%%%%%%%%%%%%%%%%%%%%%%%%%%%%%%%%%%%%%%%%%%%%%%%%%%%%%%%%

Seasonality is a repeating pattern with a known or approximately fixed
period.

Examples include:

\[
\boxed{
\text{daily cycles},
}
\]

\[
\boxed{
\text{weekly cycles},
}
\]

or

\[
\boxed{
\text{annual cycles}.
}
\]

A monthly series with yearly seasonality has period

\[
\boxed{
s=12.
}
\]

%%%%%%%%%%%%%%%%%%%%%%%%%%%%%%%%%%%%%%%%%%%%%%%%%%%%%%%%%%%%
\subsection{Cyclic Variation}
\label{subsec:time-series-cycles}
%%%%%%%%%%%%%%%%%%%%%%%%%%%%%%%%%%%%%%%%%%%%%%%%%%%%%%%%%%%%

Cyclic variation involves recurring rises and falls whose periods are not
necessarily fixed.

Economic business cycles are a common example.

Unlike strict seasonality, cycles may vary in duration and magnitude.

%%%%%%%%%%%%%%%%%%%%%%%%%%%%%%%%%%%%%%%%%%%%%%%%%%%%%%%%%%%%
\subsection{Irregular Variation}
\label{subsec:time-series-irregular}
%%%%%%%%%%%%%%%%%%%%%%%%%%%%%%%%%%%%%%%%%%%%%%%%%%%%%%%%%%%%

Irregular variation refers to remaining fluctuations not explained by
trend, seasonality, or other systematic structure.

A statistical model attempts to make this remaining component behave as
close as possible to an unpredictable noise process.

%%%%%%%%%%%%%%%%%%%%%%%%%%%%%%%%%%%%%%%%%%%%%%%%%%%%%%%%%%%%
\subsection{Additive Decomposition}
\label{subsec:additive-time-series-decomposition}
%%%%%%%%%%%%%%%%%%%%%%%%%%%%%%%%%%%%%%%%%%%%%%%%%%%%%%%%%%%%

An additive decomposition has the form

\[
\boxed{
X_t
=
T_t
+
S_t
+
C_t
+
R_t,
}
\]

where

\[
T_t
=
\text{trend},
\]

\[
S_t
=
\text{seasonal component},
\]

\[
C_t
=
\text{cyclic component},
\]

and

\[
R_t
=
\text{remainder}.
\]

%%%%%%%%%%%%%%%%%%%%%%%%%%%%%%%%%%%%%%%%%%%%%%%%%%%%%%%%%%%%
\subsection{Multiplicative Decomposition}
\label{subsec:multiplicative-time-series-decomposition}
%%%%%%%%%%%%%%%%%%%%%%%%%%%%%%%%%%%%%%%%%%%%%%%%%%%%%%%%%%%%

When variation increases with the level of the series, a multiplicative
form may be more appropriate:

\[
\boxed{
X_t
=
T_tS_tC_tR_t.
}
\]

Taking logarithms converts products into sums:

\[
\boxed{
\log X_t
=
\log T_t
+
\log S_t
+
\log C_t
+
\log R_t.
}
\]

%%%%%%%%%%%%%%%%%%%%%%%%%%%%%%%%%%%%%%%%%%%%%%%%%%%%%%%%%%%%
\subsection{Lag}
\label{subsec:time-series-lag}
%%%%%%%%%%%%%%%%%%%%%%%%%%%%%%%%%%%%%%%%%%%%%%%%%%%%%%%%%%%%

A lag is a time separation between observations.

For example,

\[
X_t
\]

and

\[
X_{t-1}
\]

are separated by lag \(1\).

Likewise,

\[
X_t
\]

and

\[
X_{t-k}
\]

are separated by lag \(k\).

%%%%%%%%%%%%%%%%%%%%%%%%%%%%%%%%%%%%%%%%%%%%%%%%%%%%%%%%%%%%
\subsection{Mean Function}
\label{subsec:time-series-mean-function}
%%%%%%%%%%%%%%%%%%%%%%%%%%%%%%%%%%%%%%%%%%%%%%%%%%%%%%%%%%%%

The time-dependent mean is

\[
\boxed{
\mu_t
=
\mathbb{E}[X_t].
}
\]

For a stationary process, this mean does not depend on time.

Thus,

\[
\boxed{
\mu_t=\mu.
}
\]

%%%%%%%%%%%%%%%%%%%%%%%%%%%%%%%%%%%%%%%%%%%%%%%%%%%%%%%%%%%%
\subsection{Autocovariance Function}
\label{subsec:autocovariance-function}
%%%%%%%%%%%%%%%%%%%%%%%%%%%%%%%%%%%%%%%%%%%%%%%%%%%%%%%%%%%%

For a general time series,

\[
\boxed{
\gamma(s,t)
=
\operatorname{Cov}(X_s,X_t).
}
\]

For a weakly stationary process, covariance depends only on the lag

\[
h=t-s.
\]

Then

\[
\boxed{
\gamma(h)
=
\operatorname{Cov}
(
X_t,
X_{t-h}
).
}
\]

%%%%%%%%%%%%%%%%%%%%%%%%%%%%%%%%%%%%%%%%%%%%%%%%%%%%%%%%%%%%
\subsection{Autocorrelation Function}
\label{subsec:autocorrelation-function}
%%%%%%%%%%%%%%%%%%%%%%%%%%%%%%%%%%%%%%%%%%%%%%%%%%%%%%%%%%%%

The autocorrelation function, abbreviated ACF, is

\[
\boxed{
\rho(h)
=
\frac{
\gamma(h)
}{
\gamma(0)
}.
}
\]

Since

\[
\gamma(0)
=
\operatorname{Var}(X_t),
\]

we have

\[
\boxed{
\rho(0)=1.
}
\]

%%%%%%%%%%%%%%%%%%%%%%%%%%%%%%%%%%%%%%%%%%%%%%%%%%%%%%%%%%%%
\subsection{Symmetry of Autocovariance}
\label{subsec:autocovariance-symmetry}
%%%%%%%%%%%%%%%%%%%%%%%%%%%%%%%%%%%%%%%%%%%%%%%%%%%%%%%%%%%%

For a weakly stationary process,

\[
\boxed{
\gamma(-h)
=
\gamma(h).
}
\]

Therefore,

\[
\boxed{
\rho(-h)
=
\rho(h).
}
\]

Only nonnegative lags are usually displayed.

%%%%%%%%%%%%%%%%%%%%%%%%%%%%%%%%%%%%%%%%%%%%%%%%%%%%%%%%%%%%
\subsection{Weak Stationarity}
\label{subsec:weak-stationarity}
%%%%%%%%%%%%%%%%%%%%%%%%%%%%%%%%%%%%%%%%%%%%%%%%%%%%%%%%%%%%

\begin{definitionbox}{Weak Stationarity}
A process \(\{X_t\}\) is weakly stationary if:

\[
\boxed{
\mathbb{E}[X_t]
=
\mu
}
\]

for all \(t\),

\[
\boxed{
\operatorname{Var}(X_t)
=
\gamma(0)
<
\infty,
}
\]

and

\[
\boxed{
\operatorname{Cov}
(
X_t,
X_{t-h}
)
=
\gamma(h)
}
\]

depends only on lag \(h\), not on \(t\).
\end{definitionbox}

%%%%%%%%%%%%%%%%%%%%%%%%%%%%%%%%%%%%%%%%%%%%%%%%%%%%%%%%%%%%
\subsection{Strict Stationarity}
\label{subsec:strict-stationarity}
%%%%%%%%%%%%%%%%%%%%%%%%%%%%%%%%%%%%%%%%%%%%%%%%%%%%%%%%%%%%

A process is strictly stationary if every finite-dimensional
distribution is unchanged by shifting time.

For any times

\[
t_1,\ldots,t_k
\]

and shift \(h\),

\[
\boxed{
(
X_{t_1},
\ldots,
X_{t_k}
)
\overset{d}{=}
(
X_{t_1+h},
\ldots,
X_{t_k+h}
).
}
\]

Strict stationarity is stronger than weak stationarity unless additional
conditions are imposed.

%%%%%%%%%%%%%%%%%%%%%%%%%%%%%%%%%%%%%%%%%%%%%%%%%%%%%%%%%%%%
\subsection{Gaussian Stationarity}
\label{subsec:gaussian-stationarity}
%%%%%%%%%%%%%%%%%%%%%%%%%%%%%%%%%%%%%%%%%%%%%%%%%%%%%%%%%%%%

For a Gaussian process, the joint distributions are completely
determined by means and covariances.

Therefore, for Gaussian processes,

\[
\boxed{
\text{weak stationarity}
\Longrightarrow
\text{strict stationarity}.
}
\]

This is a special property of Gaussian processes.

%%%%%%%%%%%%%%%%%%%%%%%%%%%%%%%%%%%%%%%%%%%%%%%%%%%%%%%%%%%%
\subsection{Sample Mean of a Time Series}
\label{subsec:time-series-sample-mean}
%%%%%%%%%%%%%%%%%%%%%%%%%%%%%%%%%%%%%%%%%%%%%%%%%%%%%%%%%%%%

For observations

\[
X_1,\ldots,X_T,
\]

the sample mean is

\[
\boxed{
\overline{X}
=
\frac1T
\sum_{t=1}^{T}
X_t.
}
\]

Because observations may be dependent, its variance is generally not

\[
\sigma^2/T.
\]

Autocovariances contribute to its uncertainty.

%%%%%%%%%%%%%%%%%%%%%%%%%%%%%%%%%%%%%%%%%%%%%%%%%%%%%%%%%%%%
\subsection{Variance of a Time-Series Mean}
\label{subsec:variance-time-series-mean}
%%%%%%%%%%%%%%%%%%%%%%%%%%%%%%%%%%%%%%%%%%%%%%%%%%%%%%%%%%%%

For a weakly stationary process,

\[
\operatorname{Var}
(
\overline{X}
)
=
\frac1{T^2}
\sum_{t=1}^{T}
\sum_{s=1}^{T}
\gamma(t-s).
\]

This can be rewritten as

\[
\boxed{
\operatorname{Var}
(
\overline{X}
)
=
\frac1T
\left[
\gamma(0)
+
2
\sum_{h=1}^{T-1}
\left(
1-\frac{h}{T}
\right)
\gamma(h)
\right].
}
\]

Positive serial correlation generally increases uncertainty in the sample
mean.

%%%%%%%%%%%%%%%%%%%%%%%%%%%%%%%%%%%%%%%%%%%%%%%%%%%%%%%%%%%%
\subsection{Sample Autocovariance}
\label{subsec:sample-autocovariance}
%%%%%%%%%%%%%%%%%%%%%%%%%%%%%%%%%%%%%%%%%%%%%%%%%%%%%%%%%%%%

A common sample autocovariance estimator at lag \(h\) is

\[
\boxed{
\widehat{\gamma}(h)
=
\frac1T
\sum_{t=h+1}^{T}
(
X_t-\overline{X}
)
(
X_{t-h}-\overline{X}
).
}
\]

Different denominator conventions are used.

%%%%%%%%%%%%%%%%%%%%%%%%%%%%%%%%%%%%%%%%%%%%%%%%%%%%%%%%%%%%
\subsection{Sample Autocorrelation}
\label{subsec:sample-autocorrelation}
%%%%%%%%%%%%%%%%%%%%%%%%%%%%%%%%%%%%%%%%%%%%%%%%%%%%%%%%%%%%

The sample autocorrelation at lag \(h\) is

\[
\boxed{
\widehat{\rho}(h)
=
\frac{
\widehat{\gamma}(h)
}{
\widehat{\gamma}(0)
}.
}
\]

Plotting

\[
\widehat{\rho}(h)
\]

against \(h\) produces the sample ACF.

%%%%%%%%%%%%%%%%%%%%%%%%%%%%%%%%%%%%%%%%%%%%%%%%%%%%%%%%%%%%
\subsection{Autocorrelation Plot}
\label{subsec:acf-plot}
%%%%%%%%%%%%%%%%%%%%%%%%%%%%%%%%%%%%%%%%%%%%%%%%%%%%%%%%%%%%

An ACF plot helps identify:

\[
\boxed{
\text{serial dependence},
}
\]

\[
\boxed{
\text{seasonality},
}
\]

and

\[
\boxed{
\text{possible ARMA structure}.
}
\]

Large autocorrelations at seasonal lags may indicate repeating seasonal
behavior.

%%%%%%%%%%%%%%%%%%%%%%%%%%%%%%%%%%%%%%%%%%%%%%%%%%%%%%%%%%%%
\subsection{White Noise}
\label{subsec:white-noise}
%%%%%%%%%%%%%%%%%%%%%%%%%%%%%%%%%%%%%%%%%%%%%%%%%%%%%%%%%%%%

\begin{definitionbox}{White Noise}
A white-noise process \(\{\varepsilon_t\}\) has

\[
\boxed{
\mathbb{E}[\varepsilon_t]=0,
}
\]

\[
\boxed{
\operatorname{Var}(\varepsilon_t)=\sigma_\varepsilon^2,
}
\]

and

\[
\boxed{
\operatorname{Cov}
(
\varepsilon_t,
\varepsilon_s
)
=
0
\qquad
(t\neq s).
}
\]
\end{definitionbox}

White noise is uncorrelated across time.

%%%%%%%%%%%%%%%%%%%%%%%%%%%%%%%%%%%%%%%%%%%%%%%%%%%%%%%%%%%%
\subsection{Independent White Noise}
\label{subsec:independent-white-noise}
%%%%%%%%%%%%%%%%%%%%%%%%%%%%%%%%%%%%%%%%%%%%%%%%%%%%%%%%%%%%

If the white-noise variables are also independent, the process is
stronger than merely uncorrelated white noise.

Gaussian white noise satisfies

\[
\boxed{
\varepsilon_t
\overset{\text{i.i.d.}}{\sim}
N
(
0,\sigma_\varepsilon^2
).
}
\]

In Gaussian settings, zero covariance implies independence.

%%%%%%%%%%%%%%%%%%%%%%%%%%%%%%%%%%%%%%%%%%%%%%%%%%%%%%%%%%%%
\subsection{Random Walk}
\label{subsec:random-walk}
%%%%%%%%%%%%%%%%%%%%%%%%%%%%%%%%%%%%%%%%%%%%%%%%%%%%%%%%%%%%

A random walk satisfies

\[
\boxed{
X_t
=
X_{t-1}
+
\varepsilon_t,
}
\]

where \(\varepsilon_t\) is white noise.

Iterating,

\[
\boxed{
X_t
=
X_0
+
\sum_{j=1}^{t}
\varepsilon_j.
}
\]

%%%%%%%%%%%%%%%%%%%%%%%%%%%%%%%%%%%%%%%%%%%%%%%%%%%%%%%%%%%%
\subsection{Random-Walk Mean and Variance}
\label{subsec:random-walk-mean-variance}
%%%%%%%%%%%%%%%%%%%%%%%%%%%%%%%%%%%%%%%%%%%%%%%%%%%%%%%%%%%%

If

\[
X_0
\]

is fixed and

\[
\mathbb{E}[\varepsilon_t]=0,
\]

then

\[
\boxed{
\mathbb{E}[X_t]
=
X_0.
}
\]

But

\[
\boxed{
\operatorname{Var}(X_t)
=
t\sigma_\varepsilon^2.
}
\]

The variance grows with time.

Therefore the random walk is not weakly stationary.

%%%%%%%%%%%%%%%%%%%%%%%%%%%%%%%%%%%%%%%%%%%%%%%%%%%%%%%%%%%%
\subsection{Random Walk with Drift}
\label{subsec:random-walk-drift}
%%%%%%%%%%%%%%%%%%%%%%%%%%%%%%%%%%%%%%%%%%%%%%%%%%%%%%%%%%%%

A random walk with drift is

\[
\boxed{
X_t
=
\delta
+
X_{t-1}
+
\varepsilon_t.
}
\]

Then

\[
\boxed{
\mathbb{E}[X_t]
=
X_0+t\delta.
}
\]

Both the mean and variance change with time.

%%%%%%%%%%%%%%%%%%%%%%%%%%%%%%%%%%%%%%%%%%%%%%%%%%%%%%%%%%%%
\subsection{Differencing}
\label{subsec:differencing}
%%%%%%%%%%%%%%%%%%%%%%%%%%%%%%%%%%%%%%%%%%%%%%%%%%%%%%%%%%%%

The first difference is

\[
\boxed{
\Delta X_t
=
X_t-X_{t-1}.
}
\]

For a random walk,

\[
X_t
=
X_{t-1}
+
\varepsilon_t,
\]

so

\[
\boxed{
\Delta X_t
=
\varepsilon_t.
}
\]

Thus differencing converts the random walk into white noise.

%%%%%%%%%%%%%%%%%%%%%%%%%%%%%%%%%%%%%%%%%%%%%%%%%%%%%%%%%%%%
\subsection{Second Differences}
\label{subsec:second-differences}
%%%%%%%%%%%%%%%%%%%%%%%%%%%%%%%%%%%%%%%%%%%%%%%%%%%%%%%%%%%%

The second difference is

\[
\boxed{
\Delta^2X_t
=
\Delta
(
\Delta X_t
).
}
\]

Expanding,

\[
\boxed{
\Delta^2X_t
=
X_t
-
2X_{t-1}
+
X_{t-2}.
}
\]

Higher-order differencing can remove polynomial-type trends, but excessive
differencing can create unnecessary dependence.

%%%%%%%%%%%%%%%%%%%%%%%%%%%%%%%%%%%%%%%%%%%%%%%%%%%%%%%%%%%%
\subsection{Seasonal Differencing}
\label{subsec:seasonal-differencing}
%%%%%%%%%%%%%%%%%%%%%%%%%%%%%%%%%%%%%%%%%%%%%%%%%%%%%%%%%%%%

For seasonal period \(s\), seasonal differencing is

\[
\boxed{
\Delta_sX_t
=
X_t-X_{t-s}.
}
\]

For monthly data with yearly seasonality,

\[
\boxed{
s=12.
}
\]

Seasonal differencing can remove repeating seasonal levels.

%%%%%%%%%%%%%%%%%%%%%%%%%%%%%%%%%%%%%%%%%%%%%%%%%%%%%%%%%%%%
\subsection{Backshift Operator}
\label{subsec:backshift-operator}
%%%%%%%%%%%%%%%%%%%%%%%%%%%%%%%%%%%%%%%%%%%%%%%%%%%%%%%%%%%%

Define the backshift operator \(B\) by

\[
\boxed{
BX_t
=
X_{t-1}.
}
\]

Then

\[
\boxed{
B^kX_t
=
X_{t-k}.
}
\]

First differencing becomes

\[
\boxed{
\Delta
=
1-B.
}
\]

Thus

\[
\boxed{
\Delta X_t
=
(1-B)X_t.
}
\]

%%%%%%%%%%%%%%%%%%%%%%%%%%%%%%%%%%%%%%%%%%%%%%%%%%%%%%%%%%%%
\subsection{Seasonal Difference Operator}
\label{subsec:seasonal-backshift}
%%%%%%%%%%%%%%%%%%%%%%%%%%%%%%%%%%%%%%%%%%%%%%%%%%%%%%%%%%%%

Seasonal differencing can be written

\[
\boxed{
\Delta_s
=
1-B^s.
}
\]

Therefore,

\[
\boxed{
\Delta_sX_t
=
(
1-B^s
)
X_t.
}
\]

Operator notation makes ARIMA models compact.

%%%%%%%%%%%%%%%%%%%%%%%%%%%%%%%%%%%%%%%%%%%%%%%%%%%%%%%%%%%%
\subsection{Autoregressive Models}
\label{subsec:autoregressive-models}
%%%%%%%%%%%%%%%%%%%%%%%%%%%%%%%%%%%%%%%%%%%%%%%%%%%%%%%%%%%%

An autoregressive model expresses the current value as a linear function
of previous values plus noise.

An AR(\(p\)) model is

\[
\boxed{
X_t
=
c
+
\phi_1X_{t-1}
+
\cdots
+
\phi_pX_{t-p}
+
\varepsilon_t.
}
\]

%%%%%%%%%%%%%%%%%%%%%%%%%%%%%%%%%%%%%%%%%%%%%%%%%%%%%%%%%%%%
\subsection{AR(1) Model}
\label{subsec:ar1-model}
%%%%%%%%%%%%%%%%%%%%%%%%%%%%%%%%%%%%%%%%%%%%%%%%%%%%%%%%%%%%

The simplest autoregressive model is

\[
\boxed{
X_t
=
c
+
\phi X_{t-1}
+
\varepsilon_t.
}
\]

If

\[
|\phi|<1,
\]

the process has a stationary solution.

%%%%%%%%%%%%%%%%%%%%%%%%%%%%%%%%%%%%%%%%%%%%%%%%%%%%%%%%%%%%
\subsection{Mean of a Stationary AR(1)}
\label{subsec:ar1-mean}
%%%%%%%%%%%%%%%%%%%%%%%%%%%%%%%%%%%%%%%%%%%%%%%%%%%%%%%%%%%%

Suppose

\[
X_t
=
c+\phi X_{t-1}+\varepsilon_t
\]

is stationary.

Taking expectations,

\[
\mu
=
c+\phi\mu.
\]

Thus,

\[
\boxed{
\mu
=
\frac{
c
}{
1-\phi
}.
}
\]

%%%%%%%%%%%%%%%%%%%%%%%%%%%%%%%%%%%%%%%%%%%%%%%%%%%%%%%%%%%%
\subsection{Centered AR(1) Model}
\label{subsec:centered-ar1}
%%%%%%%%%%%%%%%%%%%%%%%%%%%%%%%%%%%%%%%%%%%%%%%%%%%%%%%%%%%%

Subtracting the stationary mean gives

\[
\boxed{
X_t-\mu
=
\phi
(
X_{t-1}-\mu
)
+
\varepsilon_t.
}
\]

This form simplifies covariance calculations.

%%%%%%%%%%%%%%%%%%%%%%%%%%%%%%%%%%%%%%%%%%%%%%%%%%%%%%%%%%%%
\subsection{Variance of a Stationary AR(1)}
\label{subsec:ar1-variance}
%%%%%%%%%%%%%%%%%%%%%%%%%%%%%%%%%%%%%%%%%%%%%%%%%%%%%%%%%%%%

For the centered model,

\[
X_t
=
\phi X_{t-1}
+
\varepsilon_t,
\]

assuming mean zero for simplicity.

Taking variances,

\[
\gamma(0)
=
\phi^2\gamma(0)
+
\sigma_\varepsilon^2.
\]

Therefore,

\[
\boxed{
\gamma(0)
=
\frac{
\sigma_\varepsilon^2
}{
1-\phi^2
}.
}
\]

%%%%%%%%%%%%%%%%%%%%%%%%%%%%%%%%%%%%%%%%%%%%%%%%%%%%%%%%%%%%
\subsection{AR(1) Autocovariance}
\label{subsec:ar1-autocovariance}
%%%%%%%%%%%%%%%%%%%%%%%%%%%%%%%%%%%%%%%%%%%%%%%%%%%%%%%%%%%%

For \(h\geq1\),

\[
\gamma(h)
=
\operatorname{Cov}
(
X_t,
X_{t-h}
).
\]

Using

\[
X_t
=
\phi X_{t-1}
+
\varepsilon_t,
\]

and independence of future innovations from past observations,

\[
\boxed{
\gamma(h)
=
\phi\gamma(h-1).
}
\]

Thus,

\[
\boxed{
\gamma(h)
=
\phi^h\gamma(0).
}
\]

%%%%%%%%%%%%%%%%%%%%%%%%%%%%%%%%%%%%%%%%%%%%%%%%%%%%%%%%%%%%
\subsection{AR(1) Autocorrelation}
\label{subsec:ar1-autocorrelation}
%%%%%%%%%%%%%%%%%%%%%%%%%%%%%%%%%%%%%%%%%%%%%%%%%%%%%%%%%%%%

Dividing by \(\gamma(0)\),

\[
\boxed{
\rho(h)
=
\phi^{|h|}.
}
\]

Therefore the ACF of an AR(1) process decays geometrically.

If

\[
0<\phi<1,
\]

the decay is positive.

If

\[
-1<\phi<0,
\]

the autocorrelations alternate in sign.

%%%%%%%%%%%%%%%%%%%%%%%%%%%%%%%%%%%%%%%%%%%%%%%%%%%%%%%%%%%%
\subsection{Worked Example: AR(1) Variance and ACF}
\label{subsec:worked-ar1}
%%%%%%%%%%%%%%%%%%%%%%%%%%%%%%%%%%%%%%%%%%%%%%%%%%%%%%%%%%%%

\begin{workedexamplebox}{Stationary AR(1)}

Suppose

\[
X_t
=
0.6X_{t-1}
+
\varepsilon_t,
\]

where

\[
\operatorname{Var}(\varepsilon_t)=4.
\]

The stationary variance is

\[
\gamma(0)
=
\frac{
4
}{
1-0.6^2
}.
\]

Thus,

\[
\begin{answerbox}
\[
\boxed{
\gamma(0)
=
6.25.
}
\]

Also,

\[
\boxed{
\rho(h)
=
0.6^{|h|}.
}
\]
\end{answerbox}

\end{workedexamplebox}

%%%%%%%%%%%%%%%%%%%%%%%%%%%%%%%%%%%%%%%%%%%%%%%%%%%%%%%%%%%%
\subsection{Infinite Moving-Average Representation of AR(1)}
\label{subsec:ar1-ma-infinity}
%%%%%%%%%%%%%%%%%%%%%%%%%%%%%%%%%%%%%%%%%%%%%%%%%%%%%%%%%%%%

If

\[
|\phi|<1,
\]

then repeatedly substituting gives

\[
X_t
=
\varepsilon_t
+
\phi\varepsilon_{t-1}
+
\phi^2\varepsilon_{t-2}
+
\cdots
\]

for the zero-mean AR(1) model.

Thus,

\[
\boxed{
X_t
=
\sum_{j=0}^{\infty}
\phi^j
\varepsilon_{t-j}.
}
\]

This representation explains stationarity when the coefficients decay
sufficiently fast.

%%%%%%%%%%%%%%%%%%%%%%%%%%%%%%%%%%%%%%%%%%%%%%%%%%%%%%%%%%%%
\subsection{AR(\(p\)) Operator Form}
\label{subsec:arp-operator-form}
%%%%%%%%%%%%%%%%%%%%%%%%%%%%%%%%%%%%%%%%%%%%%%%%%%%%%%%%%%%%

An AR(\(p\)) model may be written

\[
\boxed{
\phi(B)X_t
=
c+\varepsilon_t,
}
\]

where

\[
\boxed{
\phi(B)
=
1
-
\phi_1B
-
\cdots
-
\phi_pB^p.
}
\]

Stationarity depends on the roots of the autoregressive polynomial.

%%%%%%%%%%%%%%%%%%%%%%%%%%%%%%%%%%%%%%%%%%%%%%%%%%%%%%%%%%%%
\subsection{AR Stationarity Condition}
\label{subsec:ar-stationarity-condition}
%%%%%%%%%%%%%%%%%%%%%%%%%%%%%%%%%%%%%%%%%%%%%%%%%%%%%%%%%%%%

The AR(\(p\)) process is stationary when the roots of

\[
\boxed{
1
-
\phi_1z
-
\cdots
-
\phi_pz^p
=
0
}
\]

lie outside the unit circle:

\[
\boxed{
|z|>1.
}
\]

For AR(1), this reduces to

\[
\boxed{
|\phi|<1.
}
\]

%%%%%%%%%%%%%%%%%%%%%%%%%%%%%%%%%%%%%%%%%%%%%%%%%%%%%%%%%%%%
\subsection{Yule--Walker Equations}
\label{subsec:yule-walker}
%%%%%%%%%%%%%%%%%%%%%%%%%%%%%%%%%%%%%%%%%%%%%%%%%%%%%%%%%%%%

For a zero-mean stationary AR(\(p\)) process,

\[
X_t
=
\sum_{j=1}^{p}
\phi_jX_{t-j}
+
\varepsilon_t,
\]

multiply by \(X_{t-k}\) and take expectations.

This gives

\[
\boxed{
\gamma(k)
=
\sum_{j=1}^{p}
\phi_j
\gamma(k-j),
\qquad
k\geq1.
}
\]

These are the Yule--Walker equations.

%%%%%%%%%%%%%%%%%%%%%%%%%%%%%%%%%%%%%%%%%%%%%%%%%%%%%%%%%%%%
\subsection{Yule--Walker Matrix Form}
\label{subsec:yule-walker-matrix}
%%%%%%%%%%%%%%%%%%%%%%%%%%%%%%%%%%%%%%%%%%%%%%%%%%%%%%%%%%%%

For the first \(p\) equations,

\[
\boxed{
\begin{pmatrix}
\gamma(0)
&
\gamma(1)
&
\cdots
&
\gamma(p-1)
\\
\gamma(1)
&
\gamma(0)
&
\cdots
&
\gamma(p-2)
\\
\vdots
&
\vdots
&
\ddots
&
\vdots
\\
\gamma(p-1)
&
\gamma(p-2)
&
\cdots
&
\gamma(0)
\end{pmatrix}
\begin{pmatrix}
\phi_1\\
\phi_2\\
\vdots\\
\phi_p
\end{pmatrix}
=
\begin{pmatrix}
\gamma(1)\\
\gamma(2)\\
\vdots\\
\gamma(p)
\end{pmatrix}.
}
\]

Sample autocovariances can be substituted to estimate AR parameters.

%%%%%%%%%%%%%%%%%%%%%%%%%%%%%%%%%%%%%%%%%%%%%%%%%%%%%%%%%%%%
\subsection{Moving-Average Models}
\label{subsec:moving-average-models}
%%%%%%%%%%%%%%%%%%%%%%%%%%%%%%%%%%%%%%%%%%%%%%%%%%%%%%%%%%%%

A moving-average model expresses the current observation as a finite
linear combination of present and past innovations.

An MA(\(q\)) model is

\[
\boxed{
X_t
=
\mu
+
\varepsilon_t
+
\theta_1\varepsilon_{t-1}
+
\cdots
+
\theta_q\varepsilon_{t-q}.
}
\]

%%%%%%%%%%%%%%%%%%%%%%%%%%%%%%%%%%%%%%%%%%%%%%%%%%%%%%%%%%%%
\subsection{MA(1) Model}
\label{subsec:ma1-model}
%%%%%%%%%%%%%%%%%%%%%%%%%%%%%%%%%%%%%%%%%%%%%%%%%%%%%%%%%%%%

A zero-mean MA(1) model is

\[
\boxed{
X_t
=
\varepsilon_t
+
\theta\varepsilon_{t-1}.
}
\]

Since it is a finite linear combination of white-noise terms, it is
stationary for every finite \(\theta\).

%%%%%%%%%%%%%%%%%%%%%%%%%%%%%%%%%%%%%%%%%%%%%%%%%%%%%%%%%%%%
\subsection{Variance of an MA(1)}
\label{subsec:ma1-variance}
%%%%%%%%%%%%%%%%%%%%%%%%%%%%%%%%%%%%%%%%%%%%%%%%%%%%%%%%%%%%

For

\[
X_t
=
\varepsilon_t
+
\theta\varepsilon_{t-1},
\]

the variance is

\[
\boxed{
\gamma(0)
=
\sigma_\varepsilon^2
(
1+\theta^2
).
}
\]

This follows from independence or zero covariance of innovations at
different times.

%%%%%%%%%%%%%%%%%%%%%%%%%%%%%%%%%%%%%%%%%%%%%%%%%%%%%%%%%%%%
\subsection{MA(1) Autocovariance}
\label{subsec:ma1-autocovariance}
%%%%%%%%%%%%%%%%%%%%%%%%%%%%%%%%%%%%%%%%%%%%%%%%%%%%%%%%%%%%

At lag \(1\),

\[
\boxed{
\gamma(1)
=
\theta
\sigma_\varepsilon^2.
}
\]

For

\[
|h|>1,
\]

there are no shared innovation terms, so

\[
\boxed{
\gamma(h)=0.
}
\]

%%%%%%%%%%%%%%%%%%%%%%%%%%%%%%%%%%%%%%%%%%%%%%%%%%%%%%%%%%%%
\subsection{MA(1) Autocorrelation}
\label{subsec:ma1-autocorrelation}
%%%%%%%%%%%%%%%%%%%%%%%%%%%%%%%%%%%%%%%%%%%%%%%%%%%%%%%%%%%%

Therefore,

\[
\boxed{
\rho(1)
=
\frac{
\theta
}{
1+\theta^2
},
}
\]

and

\[
\boxed{
\rho(h)=0
\qquad
|h|>1.
}
\]

Thus the MA(1) ACF cuts off after lag \(1\).

%%%%%%%%%%%%%%%%%%%%%%%%%%%%%%%%%%%%%%%%%%%%%%%%%%%%%%%%%%%%
\subsection{MA(\(q\)) ACF Pattern}
\label{subsec:maq-acf-pattern}
%%%%%%%%%%%%%%%%%%%%%%%%%%%%%%%%%%%%%%%%%%%%%%%%%%%%%%%%%%%%

For an MA(\(q\)) process,

\[
\boxed{
\rho(h)=0
\qquad
|h|>q.
}
\]

Thus the theoretical ACF cuts off after lag \(q\).

This provides a useful identification clue.

%%%%%%%%%%%%%%%%%%%%%%%%%%%%%%%%%%%%%%%%%%%%%%%%%%%%%%%%%%%%
\subsection{Invertibility}
\label{subsec:ma-invertibility}
%%%%%%%%%%%%%%%%%%%%%%%%%%%%%%%%%%%%%%%%%%%%%%%%%%%%%%%%%%%%

Invertibility allows an MA process to be expressed uniquely as an
infinite autoregressive process.

For MA(1),

\[
\boxed{
X_t
=
\varepsilon_t
+
\theta\varepsilon_{t-1}.
}
\]

The usual invertibility condition is

\[
\boxed{
|\theta|<1.
}
\]

%%%%%%%%%%%%%%%%%%%%%%%%%%%%%%%%%%%%%%%%%%%%%%%%%%%%%%%%%%%%
\subsection{General MA Invertibility}
\label{subsec:general-ma-invertibility}
%%%%%%%%%%%%%%%%%%%%%%%%%%%%%%%%%%%%%%%%%%%%%%%%%%%%%%%%%%%%

For

\[
\theta(B)
=
1
+
\theta_1B
+
\cdots
+
\theta_qB^q,
\]

invertibility requires the roots of

\[
\boxed{
\theta(z)=0
}
\]

to lie outside the unit circle.

This provides a unique and stable innovation representation.

%%%%%%%%%%%%%%%%%%%%%%%%%%%%%%%%%%%%%%%%%%%%%%%%%%%%%%%%%%%%
\subsection{ARMA Models}
\label{subsec:arma-models}
%%%%%%%%%%%%%%%%%%%%%%%%%%%%%%%%%%%%%%%%%%%%%%%%%%%%%%%%%%%%

An ARMA(\(p,q\)) model combines autoregressive and moving-average terms:

\[
\boxed{
X_t
=
c
+
\sum_{j=1}^{p}
\phi_jX_{t-j}
+
\varepsilon_t
+
\sum_{k=1}^{q}
\theta_k\varepsilon_{t-k}.
}
\]

In operator form,

\[
\boxed{
\phi(B)X_t
=
c
+
\theta(B)\varepsilon_t.
}
\]

%%%%%%%%%%%%%%%%%%%%%%%%%%%%%%%%%%%%%%%%%%%%%%%%%%%%%%%%%%%%
\subsection{ARMA Identification}
\label{subsec:arma-identification}
%%%%%%%%%%%%%%%%%%%%%%%%%%%%%%%%%%%%%%%%%%%%%%%%%%%%%%%%%%%%

Pure AR and MA models have characteristic ACF and PACF patterns.

For ARMA models, both the ACF and PACF generally decay rather than cut off
sharply.

Therefore model identification usually combines:

\[
\boxed{
\text{ACF},
}
\]

\[
\boxed{
\text{PACF},
}
\]

\[
\boxed{
\text{information criteria},
}
\]

and

\[
\boxed{
\text{residual diagnostics}.
}
\]

%%%%%%%%%%%%%%%%%%%%%%%%%%%%%%%%%%%%%%%%%%%%%%%%%%%%%%%%%%%%
\subsection{Partial Autocorrelation Function}
\label{subsec:partial-autocorrelation}
%%%%%%%%%%%%%%%%%%%%%%%%%%%%%%%%%%%%%%%%%%%%%%%%%%%%%%%%%%%%

The partial autocorrelation at lag \(h\) measures the linear association
between

\[
X_t
\]

and

\[
X_{t-h}
\]

after accounting for the intermediate values

\[
X_{t-1},
\ldots,
X_{t-h+1}.
\]

It is abbreviated PACF.

%%%%%%%%%%%%%%%%%%%%%%%%%%%%%%%%%%%%%%%%%%%%%%%%%%%%%%%%%%%%
\subsection{PACF Interpretation}
\label{subsec:pacf-interpretation}
%%%%%%%%%%%%%%%%%%%%%%%%%%%%%%%%%%%%%%%%%%%%%%%%%%%%%%%%%%%%

For an AR(\(p\)) process, the theoretical PACF cuts off after lag \(p\).

For an MA(\(q\)) process, the theoretical PACF generally decays.

This complements the ACF behavior.

%%%%%%%%%%%%%%%%%%%%%%%%%%%%%%%%%%%%%%%%%%%%%%%%%%%%%%%%%%%%
\subsection{Basic ACF and PACF Patterns}
\label{subsec:acf-pacf-patterns}
%%%%%%%%%%%%%%%%%%%%%%%%%%%%%%%%%%%%%%%%%%%%%%%%%%%%%%%%%%%%

A common identification summary is:

\[
\begin{array}{c|c|c}
\text{Model}
&
\text{ACF}
&
\text{PACF}
\\
\hline
AR(p)
&
\text{decays}
&
\text{cuts off after }p
\\
MA(q)
&
\text{cuts off after }q
&
\text{decays}
\\
ARMA(p,q)
&
\text{decays}
&
\text{decays}
\end{array}
\]

These are theoretical patterns, not exact finite-sample rules.

%%%%%%%%%%%%%%%%%%%%%%%%%%%%%%%%%%%%%%%%%%%%%%%%%%%%%%%%%%%%
\subsection{ARIMA Models}
\label{subsec:arima-models}
%%%%%%%%%%%%%%%%%%%%%%%%%%%%%%%%%%%%%%%%%%%%%%%%%%%%%%%%%%%%

ARIMA stands for

\[
\boxed{
\text{Autoregressive Integrated Moving Average}.
}
\]

An ARIMA(\(p,d,q\)) model assumes that after differencing \(d\) times,
the resulting series follows an ARMA(\(p,q\)) process.

Symbolically,

\[
\boxed{
\phi(B)
(
1-B
)^d
X_t
=
c
+
\theta(B)\varepsilon_t.
}
\]

%%%%%%%%%%%%%%%%%%%%%%%%%%%%%%%%%%%%%%%%%%%%%%%%%%%%%%%%%%%%
\subsection{Meaning of Integrated}
\label{subsec:meaning-of-integrated-arima}
%%%%%%%%%%%%%%%%%%%%%%%%%%%%%%%%%%%%%%%%%%%%%%%%%%%%%%%%%%%%

The term ``integrated'' refers to the fact that differencing reverses a
cumulative or integrated behavior.

For example, a random walk is ARIMA(\(0,1,0\)):

\[
\boxed{
(1-B)X_t
=
\varepsilon_t.
}
\]

Thus its first difference is white noise.

%%%%%%%%%%%%%%%%%%%%%%%%%%%%%%%%%%%%%%%%%%%%%%%%%%%%%%%%%%%%
\subsection{Worked Example: ARIMA Identification}
\label{subsec:worked-arima-identification}
%%%%%%%%%%%%%%%%%%%%%%%%%%%%%%%%%%%%%%%%%%%%%%%%%%%%%%%%%%%%

\begin{workedexamplebox}{Random Walk}

Suppose

\[
X_t
=
X_{t-1}
+
\varepsilon_t.
\]

Subtract \(X_{t-1}\):

\[
X_t-X_{t-1}
=
\varepsilon_t.
\]

Therefore,

\begin{answerbox}
\[
\boxed{
\Delta X_t
=
\varepsilon_t,
}
\]

so the process is

\[
\boxed{
ARIMA(0,1,0).
}
\]
\end{answerbox}

\end{workedexamplebox}

%%%%%%%%%%%%%%%%%%%%%%%%%%%%%%%%%%%%%%%%%%%%%%%%%%%%%%%%%%%%
\subsection{Seasonal ARIMA}
\label{subsec:seasonal-arima}
%%%%%%%%%%%%%%%%%%%%%%%%%%%%%%%%%%%%%%%%%%%%%%%%%%%%%%%%%%%%

A seasonal ARIMA model is commonly denoted

\[
\boxed{
ARIMA
(
p,d,q
)
\times
(
P,D,Q
)_s.
}
\]

Here:

\[
p,d,q
\]

describe nonseasonal structure,

while

\[
P,D,Q
\]

describe seasonal structure with period

\[
s.
\]

%%%%%%%%%%%%%%%%%%%%%%%%%%%%%%%%%%%%%%%%%%%%%%%%%%%%%%%%%%%%
\subsection{Seasonal ARIMA Operator Form}
\label{subsec:seasonal-arima-operator}
%%%%%%%%%%%%%%%%%%%%%%%%%%%%%%%%%%%%%%%%%%%%%%%%%%%%%%%%%%%%

A general multiplicative seasonal ARIMA model can be written

\[
\boxed{
\Phi
(
B^s
)
\phi(B)
(
1-B
)^d
(
1-B^s
)^D
X_t
=
\Theta
(
B^s
)
\theta(B)
\varepsilon_t.
}
\]

This notation compactly represents both ordinary and seasonal
dependence.

%%%%%%%%%%%%%%%%%%%%%%%%%%%%%%%%%%%%%%%%%%%%%%%%%%%%%%%%%%%%
\subsection{Deterministic Trend}
\label{subsec:deterministic-trend}
%%%%%%%%%%%%%%%%%%%%%%%%%%%%%%%%%%%%%%%%%%%%%%%%%%%%%%%%%%%%

A deterministic trend model may be written

\[
\boxed{
X_t
=
m(t)
+
Y_t,
}
\]

where

\[
m(t)
\]

is a deterministic function and

\[
Y_t
\]

is stationary.

For a linear trend,

\[
\boxed{
m(t)
=
\beta_0+\beta_1t.
}
\]

Detrending can be performed by estimating and subtracting \(m(t)\).

%%%%%%%%%%%%%%%%%%%%%%%%%%%%%%%%%%%%%%%%%%%%%%%%%%%%%%%%%%%%
\subsection{Stochastic Trend}
\label{subsec:stochastic-trend}
%%%%%%%%%%%%%%%%%%%%%%%%%%%%%%%%%%%%%%%%%%%%%%%%%%%%%%%%%%%%

A stochastic trend evolves through random shocks.

The random walk

\[
X_t
=
X_{t-1}
+
\varepsilon_t
\]

is the canonical example.

Shocks have persistent effects because each innovation remains embedded
in all future values.

Differencing is often appropriate for stochastic trends.

%%%%%%%%%%%%%%%%%%%%%%%%%%%%%%%%%%%%%%%%%%%%%%%%%%%%%%%%%%%%
\subsection{Detrending Versus Differencing}
\label{subsec:detrending-versus-differencing}
%%%%%%%%%%%%%%%%%%%%%%%%%%%%%%%%%%%%%%%%%%%%%%%%%%%%%%%%%%%%

For deterministic trends, regression on time may produce a stationary
remainder.

For stochastic trends, simply subtracting a deterministic trend may not
remove the nonstationarity.

Thus:

\[
\boxed{
\text{deterministic trend}
\rightarrow
\text{detrending},
}
\]

while

\[
\boxed{
\text{unit-root-type stochastic trend}
\rightarrow
\text{differencing}.
}
\]

%%%%%%%%%%%%%%%%%%%%%%%%%%%%%%%%%%%%%%%%%%%%%%%%%%%%%%%%%%%%
\subsection{Unit Roots}
\label{subsec:unit-roots}
%%%%%%%%%%%%%%%%%%%%%%%%%%%%%%%%%%%%%%%%%%%%%%%%%%%%%%%%%%%%

Consider

\[
X_t
=
\phi X_{t-1}
+
\varepsilon_t.
\]

When

\[
\boxed{
\phi=1,
}
\]

the autoregressive polynomial has a root on the unit circle.

The process becomes a random walk:

\[
\boxed{
X_t
=
X_{t-1}
+
\varepsilon_t.
}
\]

This is called a unit-root process.

%%%%%%%%%%%%%%%%%%%%%%%%%%%%%%%%%%%%%%%%%%%%%%%%%%%%%%%%%%%%
\subsection{Near Unit Roots}
\label{subsec:near-unit-roots}
%%%%%%%%%%%%%%%%%%%%%%%%%%%%%%%%%%%%%%%%%%%%%%%%%%%%%%%%%%%%

If

\[
\phi
\]

is close to \(1\), the process is stationary when

\[
|\phi|<1,
\]

but highly persistent.

Autocorrelations decay slowly.

Finite samples can make it difficult to distinguish a highly persistent
stationary process from a true unit-root process.

%%%%%%%%%%%%%%%%%%%%%%%%%%%%%%%%%%%%%%%%%%%%%%%%%%%%%%%%%%%%
\subsection{Unit-Root Testing Preview}
\label{subsec:unit-root-testing-preview}
%%%%%%%%%%%%%%%%%%%%%%%%%%%%%%%%%%%%%%%%%%%%%%%%%%%%%%%%%%%%

A common test considers

\[
H_0:
\phi=1
\]

against

\[
H_A:
|\phi|<1.
\]

Ordinary \(t\)-test reference distributions do not apply under the unit
root null.

Tests such as the Dickey--Fuller family use specialized nonstandard null
distributions.

%%%%%%%%%%%%%%%%%%%%%%%%%%%%%%%%%%%%%%%%%%%%%%%%%%%%%%%%%%%%
\subsection{Spurious Regression}
\label{subsec:spurious-regression}
%%%%%%%%%%%%%%%%%%%%%%%%%%%%%%%%%%%%%%%%%%%%%%%%%%%%%%%%%%%%

Regressing one unrelated nonstationary series on another can produce:

\[
\boxed{
\text{high }R^2,
}
\]

\[
\boxed{
\text{apparently significant coefficients},
}
\]

and

\[
\boxed{
\text{misleading relationships}.
}
\]

This phenomenon is called spurious regression.

Stationarity and cointegration must therefore be considered before
interpreting time-series regressions.

%%%%%%%%%%%%%%%%%%%%%%%%%%%%%%%%%%%%%%%%%%%%%%%%%%%%%%%%%%%%
\subsection{Transformations}
\label{subsec:time-series-transformations}
%%%%%%%%%%%%%%%%%%%%%%%%%%%%%%%%%%%%%%%%%%%%%%%%%%%%%%%%%%%%

Transformations can stabilize variance or simplify temporal structure.

Common transformations include:

\[
\boxed{
\log X_t,
}
\]

\[
\boxed{
\sqrt{X_t},
}
\]

or more generally

\[
\boxed{
g(X_t).
}
\]

The transformation should respect the domain and scientific meaning of
the data.

%%%%%%%%%%%%%%%%%%%%%%%%%%%%%%%%%%%%%%%%%%%%%%%%%%%%%%%%%%%%
\subsection{Log Differences}
\label{subsec:log-differences}
%%%%%%%%%%%%%%%%%%%%%%%%%%%%%%%%%%%%%%%%%%%%%%%%%%%%%%%%%%%%

For positive \(X_t\),

\[
\boxed{
\Delta\log X_t
=
\log X_t
-
\log X_{t-1}
=
\log
\left(
\frac{
X_t
}{
X_{t-1}
}
\right).
}
\]

For small relative changes,

\[
\boxed{
\Delta\log X_t
\approx
\frac{
X_t-X_{t-1}
}{
X_{t-1}
}.
}
\]

Thus log differences approximate proportional growth rates.

%%%%%%%%%%%%%%%%%%%%%%%%%%%%%%%%%%%%%%%%%%%%%%%%%%%%%%%%%%%%
\subsection{Box--Cox Transformation Preview}
\label{subsec:box-cox-time-series}
%%%%%%%%%%%%%%%%%%%%%%%%%%%%%%%%%%%%%%%%%%%%%%%%%%%%%%%%%%%%

For positive data, the Box--Cox family is

\[
\boxed{
g_\lambda(x)
=
\begin{cases}
\dfrac{x^\lambda-1}{\lambda},
&
\lambda\neq0,
\\[1em]
\log x,
&
\lambda=0.
\end{cases}
}
\]

It provides a systematic family for variance-stabilizing
transformations.

%%%%%%%%%%%%%%%%%%%%%%%%%%%%%%%%%%%%%%%%%%%%%%%%%%%%%%%%%%%%
\subsection{Forecasting}
\label{subsec:time-series-forecasting}
%%%%%%%%%%%%%%%%%%%%%%%%%%%%%%%%%%%%%%%%%%%%%%%%%%%%%%%%%%%%

A forecast predicts future observations using information available up to
the present.

Let

\[
\mathcal{F}_t
\]

represent information observed through time \(t\).

Under squared-error loss, the optimal forecast of \(X_{t+h}\) is

\[
\boxed{
\widehat{X}_{t+h\mid t}
=
\mathbb{E}
[
X_{t+h}
\mid
\mathcal{F}_t
].
}
\]

%%%%%%%%%%%%%%%%%%%%%%%%%%%%%%%%%%%%%%%%%%%%%%%%%%%%%%%%%%%%
\subsection{Forecast Error}
\label{subsec:forecast-error}
%%%%%%%%%%%%%%%%%%%%%%%%%%%%%%%%%%%%%%%%%%%%%%%%%%%%%%%%%%%%

The \(h\)-step-ahead forecast error is

\[
\boxed{
e_{t+h\mid t}
=
X_{t+h}
-
\widehat{X}_{t+h\mid t}.
}
\]

Its variance determines forecast uncertainty.

As the forecast horizon increases, uncertainty often increases.

%%%%%%%%%%%%%%%%%%%%%%%%%%%%%%%%%%%%%%%%%%%%%%%%%%%%%%%%%%%%
\subsection{One-Step Forecast for AR(1)}
\label{subsec:ar1-one-step-forecast}
%%%%%%%%%%%%%%%%%%%%%%%%%%%%%%%%%%%%%%%%%%%%%%%%%%%%%%%%%%%%

For

\[
X_t
=
c
+
\phi X_{t-1}
+
\varepsilon_t,
\]

the one-step forecast is

\[
\boxed{
\widehat{X}_{t+1\mid t}
=
c
+
\phi X_t.
}
\]

The forecast error is

\[
\boxed{
e_{t+1\mid t}
=
\varepsilon_{t+1}.
}
\]

Therefore,

\[
\boxed{
\operatorname{Var}
(
e_{t+1\mid t}
)
=
\sigma_\varepsilon^2.
}
\]

%%%%%%%%%%%%%%%%%%%%%%%%%%%%%%%%%%%%%%%%%%%%%%%%%%%%%%%%%%%%
\subsection{Multi-Step AR(1) Forecast}
\label{subsec:ar1-multistep-forecast}
%%%%%%%%%%%%%%%%%%%%%%%%%%%%%%%%%%%%%%%%%%%%%%%%%%%%%%%%%%%%

For a stationary AR(1) with mean \(\mu\),

\[
X_t-\mu
=
\phi
(
X_{t-1}-\mu
)
+
\varepsilon_t.
\]

The \(h\)-step forecast is

\[
\boxed{
\widehat{X}_{t+h\mid t}
=
\mu
+
\phi^h
(
X_t-\mu
).
}
\]

As

\[
h\to\infty,
\]

if \(|\phi|<1\),

\[
\boxed{
\widehat{X}_{t+h\mid t}
\to
\mu.
}
\]

%%%%%%%%%%%%%%%%%%%%%%%%%%%%%%%%%%%%%%%%%%%%%%%%%%%%%%%%%%%%
\subsection{AR(1) Forecast Error Variance}
\label{subsec:ar1-forecast-error-variance}
%%%%%%%%%%%%%%%%%%%%%%%%%%%%%%%%%%%%%%%%%%%%%%%%%%%%%%%%%%%%

For the stationary AR(1),

\[
e_{t+h\mid t}
=
\sum_{j=0}^{h-1}
\phi^j
\varepsilon_{t+h-j}.
\]

Thus,

\[
\boxed{
\operatorname{Var}
(
e_{t+h\mid t}
)
=
\sigma_\varepsilon^2
\sum_{j=0}^{h-1}
\phi^{2j}.
}
\]

Therefore,

\[
\boxed{
\operatorname{Var}
(
e_{t+h\mid t}
)
=
\sigma_\varepsilon^2
\frac{
1-\phi^{2h}
}{
1-\phi^2
}.
}
\]

%%%%%%%%%%%%%%%%%%%%%%%%%%%%%%%%%%%%%%%%%%%%%%%%%%%%%%%%%%%%
\subsection{Long-Horizon AR(1) Forecast Uncertainty}
\label{subsec:ar1-long-horizon-uncertainty}
%%%%%%%%%%%%%%%%%%%%%%%%%%%%%%%%%%%%%%%%%%%%%%%%%%%%%%%%%%%%

As

\[
h\to\infty,
\]

\[
\phi^{2h}\to0
\]

for \(|\phi|<1\).

Therefore,

\[
\boxed{
\operatorname{Var}
(
e_{t+h\mid t}
)
\to
\frac{
\sigma_\varepsilon^2
}{
1-\phi^2
}.
}
\]

This equals the unconditional stationary variance.

At long horizons, current information becomes less useful.

%%%%%%%%%%%%%%%%%%%%%%%%%%%%%%%%%%%%%%%%%%%%%%%%%%%%%%%%%%%%
\subsection{Forecast Intervals}
\label{subsec:forecast-intervals-time-series}
%%%%%%%%%%%%%%%%%%%%%%%%%%%%%%%%%%%%%%%%%%%%%%%%%%%%%%%%%%%%

If forecast errors are approximately normal,

\[
\boxed{
\widehat{X}_{t+h\mid t}
\pm
z_{1-\alpha/2}
\sqrt{
\operatorname{Var}
(
e_{t+h\mid t}
)
}
}
\]

gives an approximate forecast interval when model parameters are treated
as known.

In practice, parameter estimation adds additional uncertainty.

%%%%%%%%%%%%%%%%%%%%%%%%%%%%%%%%%%%%%%%%%%%%%%%%%%%%%%%%%%%%
\subsection{Forecasting a Random Walk}
\label{subsec:random-walk-forecast}
%%%%%%%%%%%%%%%%%%%%%%%%%%%%%%%%%%%%%%%%%%%%%%%%%%%%%%%%%%%%

For

\[
X_t
=
X_{t-1}
+
\varepsilon_t,
\]

the optimal forecast is

\[
\boxed{
\widehat{X}_{t+h\mid t}
=
X_t.
}
\]

The future value is the current value plus future innovations.

The forecast error variance is

\[
\boxed{
h\sigma_\varepsilon^2.
}
\]

Thus forecast uncertainty grows without bound.

%%%%%%%%%%%%%%%%%%%%%%%%%%%%%%%%%%%%%%%%%%%%%%%%%%%%%%%%%%%%
\subsection{Forecast Accuracy}
\label{subsec:forecast-accuracy}
%%%%%%%%%%%%%%%%%%%%%%%%%%%%%%%%%%%%%%%%%%%%%%%%%%%%%%%%%%%%

Common forecast error measures include mean absolute error:

\[
\boxed{
MAE
=
\frac1n
\sum_{i=1}^{n}
|e_i|,
}
\]

and root mean squared error:

\[
\boxed{
RMSE
=
\sqrt{
\frac1n
\sum_{i=1}^{n}
e_i^2
}.
}
\]

Different measures penalize errors differently.

%%%%%%%%%%%%%%%%%%%%%%%%%%%%%%%%%%%%%%%%%%%%%%%%%%%%%%%%%%%%
\subsection{Mean Absolute Percentage Error}
\label{subsec:mape}
%%%%%%%%%%%%%%%%%%%%%%%%%%%%%%%%%%%%%%%%%%%%%%%%%%%%%%%%%%%%

A common scale-free measure is

\[
\boxed{
MAPE
=
\frac{100}{n}
\sum_{i=1}^{n}
\left|
\frac{
e_i
}{
y_i
}
\right|.
}
\]

MAPE can behave poorly when actual values are zero or close to zero.

Thus it should not be used mechanically.

%%%%%%%%%%%%%%%%%%%%%%%%%%%%%%%%%%%%%%%%%%%%%%%%%%%%%%%%%%%%
\subsection{Training and Test Periods}
\label{subsec:time-series-training-test}
%%%%%%%%%%%%%%%%%%%%%%%%%%%%%%%%%%%%%%%%%%%%%%%%%%%%%%%%%%%%

Forecast performance should be evaluated on future observations not used
to estimate the model.

Because time order matters, randomly shuffling observations is usually
inappropriate.

Instead, use an earlier training period and a later test period:

\[
\boxed{
\text{past}
\rightarrow
\text{fit},
}
\]

\[
\boxed{
\text{future}
\rightarrow
\text{evaluate}.
}
\]

%%%%%%%%%%%%%%%%%%%%%%%%%%%%%%%%%%%%%%%%%%%%%%%%%%%%%%%%%%%%
\subsection{Rolling-Origin Evaluation}
\label{subsec:rolling-origin-evaluation}
%%%%%%%%%%%%%%%%%%%%%%%%%%%%%%%%%%%%%%%%%%%%%%%%%%%%%%%%%%%%

Rolling-origin evaluation repeatedly forecasts future observations from
different historical cutoffs.

For example:

\[
1{:}100
\rightarrow
101,
\]

\[
1{:}101
\rightarrow
102,
\]

and so forth.

This better reflects the sequential forecasting problem than random
cross-validation.

%%%%%%%%%%%%%%%%%%%%%%%%%%%%%%%%%%%%%%%%%%%%%%%%%%%%%%%%%%%%
\subsection{Expanding and Rolling Windows}
\label{subsec:expanding-rolling-windows}
%%%%%%%%%%%%%%%%%%%%%%%%%%%%%%%%%%%%%%%%%%%%%%%%%%%%%%%%%%%%

An expanding window keeps all previous observations:

\[
\boxed{
1{:}t.
}
\]

A rolling window uses only the most recent \(w\) observations:

\[
\boxed{
(t-w+1){:}t.
}
\]

Rolling windows can adapt more quickly when the data-generating process
changes over time.

%%%%%%%%%%%%%%%%%%%%%%%%%%%%%%%%%%%%%%%%%%%%%%%%%%%%%%%%%%%%
\subsection{Simple Exponential Smoothing}
\label{subsec:simple-exponential-smoothing}
%%%%%%%%%%%%%%%%%%%%%%%%%%%%%%%%%%%%%%%%%%%%%%%%%%%%%%%%%%%%

For a series with approximately constant level and no strong trend or
seasonality, simple exponential smoothing updates the level as

\[
\boxed{
\ell_t
=
\alpha X_t
+
(1-\alpha)\ell_{t-1},
}
\]

where

\[
0\leq\alpha\leq1.
\]

%%%%%%%%%%%%%%%%%%%%%%%%%%%%%%%%%%%%%%%%%%%%%%%%%%%%%%%%%%%%
\subsection{Exponential Weighting}
\label{subsec:exponential-weighting}
%%%%%%%%%%%%%%%%%%%%%%%%%%%%%%%%%%%%%%%%%%%%%%%%%%%%%%%%%%%%

Repeated substitution yields weights that decay geometrically into the
past.

Recent observations receive weights such as

\[
\alpha,
\]

\[
\alpha(1-\alpha),
\]

\[
\alpha(1-\alpha)^2,
\]

and so forth.

Thus older observations receive exponentially smaller weights.

%%%%%%%%%%%%%%%%%%%%%%%%%%%%%%%%%%%%%%%%%%%%%%%%%%%%%%%%%%%%
\subsection{Simple Exponential-Smoothing Forecast}
\label{subsec:ses-forecast}
%%%%%%%%%%%%%%%%%%%%%%%%%%%%%%%%%%%%%%%%%%%%%%%%%%%%%%%%%%%%

The forecast for all future horizons is

\[
\boxed{
\widehat{X}_{t+h\mid t}
=
\ell_t,
\qquad
h\geq1.
}
\]

Thus the model forecasts a constant future level.

%%%%%%%%%%%%%%%%%%%%%%%%%%%%%%%%%%%%%%%%%%%%%%%%%%%%%%%%%%%%
\subsection{Holt's Linear Trend Method}
\label{subsec:holt-linear-trend}
%%%%%%%%%%%%%%%%%%%%%%%%%%%%%%%%%%%%%%%%%%%%%%%%%%%%%%%%%%%%

Holt's method extends exponential smoothing by maintaining both a level
and a trend.

A common form is

\[
\boxed{
\ell_t
=
\alpha X_t
+
(1-\alpha)
(
\ell_{t-1}+b_{t-1}
),
}
\]

\[
\boxed{
b_t
=
\beta
(
\ell_t-\ell_{t-1}
)
+
(1-\beta)b_{t-1}.
}
\]

%%%%%%%%%%%%%%%%%%%%%%%%%%%%%%%%%%%%%%%%%%%%%%%%%%%%%%%%%%%%
\subsection{Holt Forecast}
\label{subsec:holt-forecast}
%%%%%%%%%%%%%%%%%%%%%%%%%%%%%%%%%%%%%%%%%%%%%%%%%%%%%%%%%%%%

The \(h\)-step forecast is

\[
\boxed{
\widehat{X}_{t+h\mid t}
=
\ell_t
+
hb_t.
}
\]

Thus the forecast extrapolates the estimated local trend.

%%%%%%%%%%%%%%%%%%%%%%%%%%%%%%%%%%%%%%%%%%%%%%%%%%%%%%%%%%%%
\subsection{Damped Trend Preview}
\label{subsec:damped-trend}
%%%%%%%%%%%%%%%%%%%%%%%%%%%%%%%%%%%%%%%%%%%%%%%%%%%%%%%%%%%%

Long-run extrapolation of a linear trend may be unrealistic.

A damped-trend method uses a damping parameter

\[
0<\phi<1
\]

so trend contributions gradually decrease with forecast horizon.

This often produces more conservative long-term forecasts.

%%%%%%%%%%%%%%%%%%%%%%%%%%%%%%%%%%%%%%%%%%%%%%%%%%%%%%%%%%%%
\subsection{Holt--Winters Methods}
\label{subsec:holt-winters}
%%%%%%%%%%%%%%%%%%%%%%%%%%%%%%%%%%%%%%%%%%%%%%%%%%%%%%%%%%%%

Holt--Winters methods extend exponential smoothing to include seasonal
components.

They maintain estimates of:

\[
\boxed{
\text{level},
}
\]

\[
\boxed{
\text{trend},
}
\]

and

\[
\boxed{
\text{seasonality}.
}
\]

Additive and multiplicative seasonal versions are commonly used.

%%%%%%%%%%%%%%%%%%%%%%%%%%%%%%%%%%%%%%%%%%%%%%%%%%%%%%%%%%%%
\subsection{Additive Versus Multiplicative Seasonality}
\label{subsec:additive-multiplicative-seasonality}
%%%%%%%%%%%%%%%%%%%%%%%%%%%%%%%%%%%%%%%%%%%%%%%%%%%%%%%%%%%%

Additive seasonality is appropriate when seasonal fluctuations have
roughly constant magnitude.

Multiplicative seasonality is appropriate when seasonal variation grows
with the level of the series.

A log transformation can often convert multiplicative structure into an
approximately additive form.

%%%%%%%%%%%%%%%%%%%%%%%%%%%%%%%%%%%%%%%%%%%%%%%%%%%%%%%%%%%%
\subsection{Time-Series Regression}
\label{subsec:time-series-regression}
%%%%%%%%%%%%%%%%%%%%%%%%%%%%%%%%%%%%%%%%%%%%%%%%%%%%%%%%%%%%

A time-series regression may have the form

\[
\boxed{
Y_t
=
\beta_0
+
\beta_1X_{1t}
+
\cdots
+
\beta_pX_{pt}
+
\varepsilon_t.
}
\]

Unlike ordinary regression, the error sequence

\[
\varepsilon_t
\]

may be serially correlated.

Ignoring that dependence can produce incorrect standard errors and
forecast intervals.

%%%%%%%%%%%%%%%%%%%%%%%%%%%%%%%%%%%%%%%%%%%%%%%%%%%%%%%%%%%%
\subsection{Regression with AR Errors}
\label{subsec:regression-ar-errors}
%%%%%%%%%%%%%%%%%%%%%%%%%%%%%%%%%%%%%%%%%%%%%%%%%%%%%%%%%%%%

One model is

\[
\boxed{
Y_t
=
\mathbf{x}_t^T
\boldsymbol{\beta}
+
u_t,
}
\]

with

\[
\boxed{
u_t
=
\phi u_{t-1}
+
\varepsilon_t.
}
\]

This combines regression on explanatory variables with autoregressive
error structure.

%%%%%%%%%%%%%%%%%%%%%%%%%%%%%%%%%%%%%%%%%%%%%%%%%%%%%%%%%%%%
\subsection{Dynamic Regression}
\label{subsec:dynamic-regression}
%%%%%%%%%%%%%%%%%%%%%%%%%%%%%%%%%%%%%%%%%%%%%%%%%%%%%%%%%%%%

A dynamic regression may include current and lagged predictors:

\[
\boxed{
Y_t
=
\beta_0
+
\beta_1X_t
+
\beta_2X_{t-1}
+
\phi Y_{t-1}
+
\varepsilon_t.
}
\]

Such models combine regression and time dependence.

Care is required when variables are nonstationary.

%%%%%%%%%%%%%%%%%%%%%%%%%%%%%%%%%%%%%%%%%%%%%%%%%%%%%%%%%%%%
\subsection{Intervention Analysis Preview}
\label{subsec:intervention-analysis}
%%%%%%%%%%%%%%%%%%%%%%%%%%%%%%%%%%%%%%%%%%%%%%%%%%%%%%%%%%%%

An intervention represents a known event that may alter a series.

For example, define

\[
I_t
=
\begin{cases}
0,&t<T_0,\\
1,&t\geq T_0.
\end{cases}
\]

A model such as

\[
\boxed{
X_t
=
\mu
+
\beta I_t
+
u_t
}
\]

can represent a permanent level change while \(u_t\) follows a time-series
model.

%%%%%%%%%%%%%%%%%%%%%%%%%%%%%%%%%%%%%%%%%%%%%%%%%%%%%%%%%%%%
\subsection{Pulse Intervention}
\label{subsec:pulse-intervention}
%%%%%%%%%%%%%%%%%%%%%%%%%%%%%%%%%%%%%%%%%%%%%%%%%%%%%%%%%%%%

A pulse intervention represents an effect occurring at one particular
time:

\[
\boxed{
P_t
=
\begin{cases}
1,&t=T_0,\\
0,&t\neq T_0.
\end{cases}
}
\]

It can represent unusual one-time events.

%%%%%%%%%%%%%%%%%%%%%%%%%%%%%%%%%%%%%%%%%%%%%%%%%%%%%%%%%%%%
\subsection{Transfer-Function Models Preview}
\label{subsec:transfer-function-models}
%%%%%%%%%%%%%%%%%%%%%%%%%%%%%%%%%%%%%%%%%%%%%%%%%%%%%%%%%%%%

A transfer-function model describes how an input series affects an output
series over time.

One may write schematically

\[
\boxed{
Y_t
=
\sum_{j=0}^{\infty}
\omega_j
X_{t-j}
+
N_t,
}
\]

where

\[
N_t
\]

is a noise process.

The coefficients \(\omega_j\) describe delayed dynamic effects.

%%%%%%%%%%%%%%%%%%%%%%%%%%%%%%%%%%%%%%%%%%%%%%%%%%%%%%%%%%%%
\subsection{Distributed-Lag Models}
\label{subsec:distributed-lag-models}
%%%%%%%%%%%%%%%%%%%%%%%%%%%%%%%%%%%%%%%%%%%%%%%%%%%%%%%%%%%%

A finite distributed-lag model is

\[
\boxed{
Y_t
=
\alpha
+
\beta_0X_t
+
\beta_1X_{t-1}
+
\cdots
+
\beta_qX_{t-q}
+
\varepsilon_t.
}
\]

The total effect of a persistent one-unit change may involve

\[
\boxed{
\sum_{j=0}^{q}\beta_j.
}
\]

%%%%%%%%%%%%%%%%%%%%%%%%%%%%%%%%%%%%%%%%%%%%%%%%%%%%%%%%%%%%
\subsection{Residuals}
\label{subsec:time-series-residuals}
%%%%%%%%%%%%%%%%%%%%%%%%%%%%%%%%%%%%%%%%%%%%%%%%%%%%%%%%%%%%

After fitting a model, residuals are

\[
\boxed{
e_t
=
X_t-\widehat{X}_t.
}
\]

A good time-series model should leave residuals that behave approximately
like unpredictable noise.

%%%%%%%%%%%%%%%%%%%%%%%%%%%%%%%%%%%%%%%%%%%%%%%%%%%%%%%%%%%%
\subsection{Residual ACF}
\label{subsec:residual-acf}
%%%%%%%%%%%%%%%%%%%%%%%%%%%%%%%%%%%%%%%%%%%%%%%%%%%%%%%%%%%%

Residual autocorrelations should be small if the temporal dependence has
been adequately modeled.

Large residual ACF values indicate remaining structure.

This may suggest:

\[
\boxed{
\text{missing AR terms},
}
\]

\[
\boxed{
\text{missing MA terms},
}
\]

\[
\boxed{
\text{seasonality},
}
\]

or

\[
\boxed{
\text{structural change}.
}
\]

%%%%%%%%%%%%%%%%%%%%%%%%%%%%%%%%%%%%%%%%%%%%%%%%%%%%%%%%%%%%
\subsection{Ljung--Box Test Preview}
\label{subsec:ljung-box-preview}
%%%%%%%%%%%%%%%%%%%%%%%%%%%%%%%%%%%%%%%%%%%%%%%%%%%%%%%%%%%%

The Ljung--Box statistic jointly evaluates several residual
autocorrelations.

For lags through \(m\),

\[
\boxed{
Q
=
T(T+2)
\sum_{h=1}^{m}
\frac{
\widehat{\rho}(h)^2
}{
T-h
}.
}
\]

Under suitable null conditions, \(Q\) has an approximate chi-square
distribution with degrees of freedom adjusted for fitted model
parameters.

%%%%%%%%%%%%%%%%%%%%%%%%%%%%%%%%%%%%%%%%%%%%%%%%%%%%%%%%%%%%
\subsection{Model Selection}
\label{subsec:time-series-model-selection}
%%%%%%%%%%%%%%%%%%%%%%%%%%%%%%%%%%%%%%%%%%%%%%%%%%%%%%%%%%%%

Time-series model selection should balance fit and complexity.

Common tools include:

\[
\boxed{
AIC,
}
\]

\[
\boxed{
BIC,
}
\]

\[
\boxed{
ACF/PACF patterns,
}
\]

\[
\boxed{
residual diagnostics,
}
\]

and

\[
\boxed{
out-of-sample forecast performance.
}
\]

No single criterion should be used blindly.

%%%%%%%%%%%%%%%%%%%%%%%%%%%%%%%%%%%%%%%%%%%%%%%%%%%%%%%%%%%%
\subsection{Akaike Information Criterion}
\label{subsec:time-series-aic}
%%%%%%%%%%%%%%%%%%%%%%%%%%%%%%%%%%%%%%%%%%%%%%%%%%%%%%%%%%%%

For a model with maximized likelihood \(L\) and \(k\) estimated
parameters,

\[
\boxed{
AIC
=
-2\log L
+
2k.
}
\]

Smaller AIC values are preferred among models fit to the same data under
comparable likelihood definitions.

%%%%%%%%%%%%%%%%%%%%%%%%%%%%%%%%%%%%%%%%%%%%%%%%%%%%%%%%%%%%
\subsection{Bayesian Information Criterion}
\label{subsec:time-series-bic}
%%%%%%%%%%%%%%%%%%%%%%%%%%%%%%%%%%%%%%%%%%%%%%%%%%%%%%%%%%%%

The Bayesian information criterion is

\[
\boxed{
BIC
=
-2\log L
+
k\log T.
}
\]

BIC penalizes model complexity more strongly than AIC when \(T\) is
sufficiently large.

%%%%%%%%%%%%%%%%%%%%%%%%%%%%%%%%%%%%%%%%%%%%%%%%%%%%%%%%%%%%
\subsection{State-Space Models}
\label{subsec:state-space-models}
%%%%%%%%%%%%%%%%%%%%%%%%%%%%%%%%%%%%%%%%%%%%%%%%%%%%%%%%%%%%

State-space models describe an observed process through an unobserved
state process.

A linear Gaussian state-space model has:

\[
\boxed{
\mathbf{x}_t
=
F\mathbf{x}_{t-1}
+
\mathbf{w}_t
}
\]

and

\[
\boxed{
\mathbf{y}_t
=
H\mathbf{x}_t
+
\mathbf{v}_t.
}
\]

The first is the state equation.

The second is the observation equation.

%%%%%%%%%%%%%%%%%%%%%%%%%%%%%%%%%%%%%%%%%%%%%%%%%%%%%%%%%%%%
\subsection{Latent States}
\label{subsec:latent-states}
%%%%%%%%%%%%%%%%%%%%%%%%%%%%%%%%%%%%%%%%%%%%%%%%%%%%%%%%%%%%

The state

\[
\mathbf{x}_t
\]

represents unobserved system information at time \(t\).

Examples include:

\[
\boxed{
\text{underlying level},
}
\]

\[
\boxed{
\text{trend},
}
\]

\[
\boxed{
\text{velocity},
}
\]

or

\[
\boxed{
\text{hidden economic condition}.
}
\]

Observed data provide noisy information about these states.

%%%%%%%%%%%%%%%%%%%%%%%%%%%%%%%%%%%%%%%%%%%%%%%%%%%%%%%%%%%%
\subsection{Kalman Filter Preview}
\label{subsec:kalman-filter-preview}
%%%%%%%%%%%%%%%%%%%%%%%%%%%%%%%%%%%%%%%%%%%%%%%%%%%%%%%%%%%%

The Kalman filter recursively estimates latent states in a linear
Gaussian state-space model.

Each step has two broad stages:

\[
\boxed{
\text{prediction}
}
\]

and

\[
\boxed{
\text{update}.
}
\]

The prediction uses the previous state estimate.

The update incorporates the new observation.

%%%%%%%%%%%%%%%%%%%%%%%%%%%%%%%%%%%%%%%%%%%%%%%%%%%%%%%%%%%%
\subsection{Kalman Prediction Step}
\label{subsec:kalman-prediction}
%%%%%%%%%%%%%%%%%%%%%%%%%%%%%%%%%%%%%%%%%%%%%%%%%%%%%%%%%%%%

If

\[
\widehat{\mathbf{x}}_{t-1\mid t-1}
\]

is the filtered state estimate, the predicted state is

\[
\boxed{
\widehat{\mathbf{x}}_{t\mid t-1}
=
F
\widehat{\mathbf{x}}_{t-1\mid t-1}.
}
\]

The covariance is also propagated through the state equation.

%%%%%%%%%%%%%%%%%%%%%%%%%%%%%%%%%%%%%%%%%%%%%%%%%%%%%%%%%%%%
\subsection{Kalman Update Idea}
\label{subsec:kalman-update}
%%%%%%%%%%%%%%%%%%%%%%%%%%%%%%%%%%%%%%%%%%%%%%%%%%%%%%%%%%%%

The observed value

\[
\mathbf{y}_t
\]

is compared with its prediction.

The difference is the innovation:

\[
\boxed{
\mathbf{v}_t
=
\mathbf{y}_t
-
H
\widehat{\mathbf{x}}_{t\mid t-1}.
}
\]

The state estimate is corrected using a matrix called the Kalman gain.

%%%%%%%%%%%%%%%%%%%%%%%%%%%%%%%%%%%%%%%%%%%%%%%%%%%%%%%%%%%%
\subsection{Structural Time-Series Models}
\label{subsec:structural-time-series}
%%%%%%%%%%%%%%%%%%%%%%%%%%%%%%%%%%%%%%%%%%%%%%%%%%%%%%%%%%%%

State-space models can represent structural components such as:

\[
\boxed{
\text{local level},
}
\]

\[
\boxed{
\text{local trend},
}
\]

\[
\boxed{
\text{seasonality},
}
\]

and

\[
\boxed{
\text{regression effects}.
}
\]

This provides an alternative framework to classical ARIMA modeling.

%%%%%%%%%%%%%%%%%%%%%%%%%%%%%%%%%%%%%%%%%%%%%%%%%%%%%%%%%%%%
\subsection{Multivariate Time Series}
\label{subsec:multivariate-time-series}
%%%%%%%%%%%%%%%%%%%%%%%%%%%%%%%%%%%%%%%%%%%%%%%%%%%%%%%%%%%%

A multivariate time series is a vector process

\[
\boxed{
\mathbf{X}_t
=
\begin{pmatrix}
X_{1t}\\
\vdots\\
X_{pt}
\end{pmatrix}.
}
\]

Now dependence can occur:

\[
\boxed{
\text{across time}
}
\]

and

\[
\boxed{
\text{across variables}.
}
\]

%%%%%%%%%%%%%%%%%%%%%%%%%%%%%%%%%%%%%%%%%%%%%%%%%%%%%%%%%%%%
\subsection{Vector Autoregression}
\label{subsec:vector-autoregression}
%%%%%%%%%%%%%%%%%%%%%%%%%%%%%%%%%%%%%%%%%%%%%%%%%%%%%%%%%%%%

A VAR(\(p\)) model is

\[
\boxed{
\mathbf{X}_t
=
\mathbf{c}
+
A_1
\mathbf{X}_{t-1}
+
\cdots
+
A_p
\mathbf{X}_{t-p}
+
\boldsymbol{\varepsilon}_t.
}
\]

Each variable may depend on past values of all variables in the system.

%%%%%%%%%%%%%%%%%%%%%%%%%%%%%%%%%%%%%%%%%%%%%%%%%%%%%%%%%%%%
\subsection{VAR(1)}
\label{subsec:var1}
%%%%%%%%%%%%%%%%%%%%%%%%%%%%%%%%%%%%%%%%%%%%%%%%%%%%%%%%%%%%

A VAR(1) model is

\[
\boxed{
\mathbf{X}_t
=
\mathbf{c}
+
A
\mathbf{X}_{t-1}
+
\boldsymbol{\varepsilon}_t.
}
\]

Its stability depends on the eigenvalues of \(A\).

A standard stationarity condition is

\[
\boxed{
|\lambda_j(A)|<1
}
\]

for every eigenvalue \(\lambda_j(A)\).

%%%%%%%%%%%%%%%%%%%%%%%%%%%%%%%%%%%%%%%%%%%%%%%%%%%%%%%%%%%%
\subsection{Cross-Covariance}
\label{subsec:cross-covariance-time-series}
%%%%%%%%%%%%%%%%%%%%%%%%%%%%%%%%%%%%%%%%%%%%%%%%%%%%%%%%%%%%

For two component series \(X_t\) and \(Y_t\), the cross-covariance at lag
\(h\) is

\[
\boxed{
\gamma_{XY}(h)
=
\operatorname{Cov}
(
X_t,
Y_{t-h}
).
}
\]

Unlike ordinary autocovariance,

\[
\gamma_{XY}(h)
\]

need not be symmetric in \(h\).

%%%%%%%%%%%%%%%%%%%%%%%%%%%%%%%%%%%%%%%%%%%%%%%%%%%%%%%%%%%%
\subsection{Granger Causality Preview}
\label{subsec:granger-causality-preview}
%%%%%%%%%%%%%%%%%%%%%%%%%%%%%%%%%%%%%%%%%%%%%%%%%%%%%%%%%%%%

A series \(X_t\) is said to Granger-cause \(Y_t\) if past values of \(X\)
improve prediction of \(Y\) beyond information already contained in past
values of \(Y\) and other included variables.

This is a predictive concept.

It does not by itself establish structural or scientific causation.

%%%%%%%%%%%%%%%%%%%%%%%%%%%%%%%%%%%%%%%%%%%%%%%%%%%%%%%%%%%%
\subsection{Cointegration}
\label{subsec:cointegration}
%%%%%%%%%%%%%%%%%%%%%%%%%%%%%%%%%%%%%%%%%%%%%%%%%%%%%%%%%%%%

Two nonstationary series may move together so that a particular linear
combination is stationary.

Suppose \(X_t\) and \(Y_t\) are individually nonstationary, but

\[
\boxed{
Y_t-\beta X_t
}
\]

is stationary.

Then the series are cointegrated.

%%%%%%%%%%%%%%%%%%%%%%%%%%%%%%%%%%%%%%%%%%%%%%%%%%%%%%%%%%%%
\subsection{Why Cointegration Matters}
\label{subsec:why-cointegration}
%%%%%%%%%%%%%%%%%%%%%%%%%%%%%%%%%%%%%%%%%%%%%%%%%%%%%%%%%%%%

Differencing removes long-run level information.

Cointegration identifies situations where nonstationary variables share a
stable long-run relationship.

Thus a model can represent:

\[
\boxed{
\text{short-run dynamics}
}
\]

and

\[
\boxed{
\text{long-run equilibrium}.
}
\]

%%%%%%%%%%%%%%%%%%%%%%%%%%%%%%%%%%%%%%%%%%%%%%%%%%%%%%%%%%%%
\subsection{Error-Correction Model Preview}
\label{subsec:error-correction-model}
%%%%%%%%%%%%%%%%%%%%%%%%%%%%%%%%%%%%%%%%%%%%%%%%%%%%%%%%%%%%

A simple error-correction model may have the form

\[
\boxed{
\Delta Y_t
=
\alpha
(
Y_{t-1}
-
\beta X_{t-1}
)
+
\gamma
\Delta X_t
+
\varepsilon_t.
}
\]

The term

\[
Y_{t-1}-\beta X_{t-1}
\]

measures deviation from the long-run relationship.

The coefficient \(\alpha\) controls adjustment toward equilibrium.

%%%%%%%%%%%%%%%%%%%%%%%%%%%%%%%%%%%%%%%%%%%%%%%%%%%%%%%%%%%%
\subsection{Frequency-Domain View}
\label{subsec:frequency-domain-view}
%%%%%%%%%%%%%%%%%%%%%%%%%%%%%%%%%%%%%%%%%%%%%%%%%%%%%%%%%%%%

Time-domain methods describe dependence through lags.

Frequency-domain methods describe a series through oscillations at
different frequencies.

Thus the same stochastic behavior can often be studied through:

\[
\boxed{
\text{time}
}
\]

or

\[
\boxed{
\text{frequency}.
}
\]

%%%%%%%%%%%%%%%%%%%%%%%%%%%%%%%%%%%%%%%%%%%%%%%%%%%%%%%%%%%%
\subsection{Sinusoidal Components}
\label{subsec:sinusoidal-components}
%%%%%%%%%%%%%%%%%%%%%%%%%%%%%%%%%%%%%%%%%%%%%%%%%%%%%%%%%%%%

A sinusoidal component has the form

\[
\boxed{
A
\cos
(
\omega t+\phi
),
}
\]

where

\[
A
=
\text{amplitude},
\]

\[
\omega
=
\text{angular frequency},
\]

and

\[
\phi
=
\text{phase}.
\]

Repeated cycles in time correspond to concentrated energy at particular
frequencies.

%%%%%%%%%%%%%%%%%%%%%%%%%%%%%%%%%%%%%%%%%%%%%%%%%%%%%%%%%%%%
\subsection{Spectral Density}
\label{subsec:spectral-density}
%%%%%%%%%%%%%%%%%%%%%%%%%%%%%%%%%%%%%%%%%%%%%%%%%%%%%%%%%%%%

For a weakly stationary process with autocovariance function
\(\gamma(h)\), the spectral density is formally

\[
\boxed{
f(\omega)
=
\frac1{2\pi}
\sum_{h=-\infty}^{\infty}
\gamma(h)
e^{-i\omega h},
}
\]

when the series converges appropriately.

It is the Fourier transform of the autocovariance function.

%%%%%%%%%%%%%%%%%%%%%%%%%%%%%%%%%%%%%%%%%%%%%%%%%%%%%%%%%%%%
\subsection{Inverse Spectral Relation}
\label{subsec:inverse-spectral-relation}
%%%%%%%%%%%%%%%%%%%%%%%%%%%%%%%%%%%%%%%%%%%%%%%%%%%%%%%%%%%%

The autocovariance can be recovered from the spectral density:

\[
\boxed{
\gamma(h)
=
\int_{-\pi}^{\pi}
e^{i\omega h}
f(\omega)
\,d\omega.
}
\]

Thus autocovariance and spectrum provide two equivalent descriptions of
stationary second-order dependence.

%%%%%%%%%%%%%%%%%%%%%%%%%%%%%%%%%%%%%%%%%%%%%%%%%%%%%%%%%%%%
\subsection{Spectrum of White Noise}
\label{subsec:white-noise-spectrum}
%%%%%%%%%%%%%%%%%%%%%%%%%%%%%%%%%%%%%%%%%%%%%%%%%%%%%%%%%%%%

For white noise,

\[
\gamma(0)
=
\sigma_\varepsilon^2
\]

and

\[
\gamma(h)=0
\qquad
h\neq0.
\]

Therefore,

\[
\boxed{
f(\omega)
=
\frac{
\sigma_\varepsilon^2
}{
2\pi
}.
}
\]

The spectrum is constant across frequencies.

This motivates the term ``white'' noise by analogy with white light.

%%%%%%%%%%%%%%%%%%%%%%%%%%%%%%%%%%%%%%%%%%%%%%%%%%%%%%%%%%%%
\subsection{Periodogram Preview}
\label{subsec:periodogram-preview}
%%%%%%%%%%%%%%%%%%%%%%%%%%%%%%%%%%%%%%%%%%%%%%%%%%%%%%%%%%%%

The periodogram estimates how sample variation is distributed across
frequencies.

A common form is

\[
\boxed{
I(\omega)
=
\frac1{2\pi T}
\left|
\sum_{t=1}^{T}
X_t
e^{-i\omega t}
\right|^2.
}
\]

Large peaks may indicate strong periodic components.

%%%%%%%%%%%%%%%%%%%%%%%%%%%%%%%%%%%%%%%%%%%%%%%%%%%%%%%%%%%%
\subsection{Time-Varying Behavior}
\label{subsec:time-varying-behavior}
%%%%%%%%%%%%%%%%%%%%%%%%%%%%%%%%%%%%%%%%%%%%%%%%%%%%%%%%%%%%

Real series may change behavior over time.

Examples include:

\[
\boxed{
\text{changing mean},
}
\]

\[
\boxed{
\text{changing variance},
}
\]

\[
\boxed{
\text{changing AR coefficients},
}
\]

or

\[
\boxed{
\text{structural breaks}.
}
\]

Stationary models may then provide only local approximations.

%%%%%%%%%%%%%%%%%%%%%%%%%%%%%%%%%%%%%%%%%%%%%%%%%%%%%%%%%%%%
\subsection{Structural Breaks}
\label{subsec:structural-breaks}
%%%%%%%%%%%%%%%%%%%%%%%%%%%%%%%%%%%%%%%%%%%%%%%%%%%%%%%%%%%%

A structural break is a persistent change in the data-generating process.

For example,

\[
\boxed{
\mu_t
=
\begin{cases}
\mu_1,&t<T_0,\\
\mu_2,&t\geq T_0.
\end{cases}
}
\]

Ignoring such a change can produce misleading estimates and forecasts.

%%%%%%%%%%%%%%%%%%%%%%%%%%%%%%%%%%%%%%%%%%%%%%%%%%%%%%%%%%%%
\subsection{Volatility Clustering Preview}
\label{subsec:volatility-clustering}
%%%%%%%%%%%%%%%%%%%%%%%%%%%%%%%%%%%%%%%%%%%%%%%%%%%%%%%%%%%%

Some time series have little autocorrelation in their levels but strong
dependence in their squared values.

Large fluctuations tend to be followed by large fluctuations.

This is called volatility clustering.

It motivates models such as ARCH and GARCH.

%%%%%%%%%%%%%%%%%%%%%%%%%%%%%%%%%%%%%%%%%%%%%%%%%%%%%%%%%%%%
\subsection{ARCH Model Preview}
\label{subsec:arch-preview}
%%%%%%%%%%%%%%%%%%%%%%%%%%%%%%%%%%%%%%%%%%%%%%%%%%%%%%%%%%%%

An ARCH(\(q\)) model allows conditional variance to depend on past squared
innovations:

\[
\boxed{
\sigma_t^2
=
\alpha_0
+
\alpha_1
\varepsilon_{t-1}^2
+
\cdots
+
\alpha_q
\varepsilon_{t-q}^2.
}
\]

Then

\[
\boxed{
\varepsilon_t
=
\sigma_tZ_t,
}
\]

where \(Z_t\) is often standardized white noise.

%%%%%%%%%%%%%%%%%%%%%%%%%%%%%%%%%%%%%%%%%%%%%%%%%%%%%%%%%%%%
\subsection{GARCH Model Preview}
\label{subsec:garch-preview}
%%%%%%%%%%%%%%%%%%%%%%%%%%%%%%%%%%%%%%%%%%%%%%%%%%%%%%%%%%%%

A GARCH(\(p,q\)) model also includes past conditional variances:

\[
\boxed{
\sigma_t^2
=
\alpha_0
+
\sum_{i=1}^{q}
\alpha_i
\varepsilon_{t-i}^2
+
\sum_{j=1}^{p}
\beta_j
\sigma_{t-j}^2.
}
\]

GARCH models are widely used for financial volatility.

%%%%%%%%%%%%%%%%%%%%%%%%%%%%%%%%%%%%%%%%%%%%%%%%%%%%%%%%%%%%
\subsection{Time-Series Modeling Workflow}
\label{subsec:time-series-workflow}
%%%%%%%%%%%%%%%%%%%%%%%%%%%%%%%%%%%%%%%%%%%%%%%%%%%%%%%%%%%%

A practical workflow is:

\begin{enumerate}
    \item plot the data against time;
    \item identify trend and seasonal structure;
    \item inspect changes in variability;
    \item transform the data if scientifically appropriate;
    \item assess stationarity;
    \item remove deterministic trend or difference as needed;
    \item inspect the ACF and PACF;
    \item identify candidate models;
    \item estimate model parameters;
    \item compare model-selection criteria;
    \item inspect residuals;
    \item inspect residual ACF;
    \item test for remaining serial dependence when useful;
    \item evaluate forecasts using future observations;
    \item quantify forecast uncertainty;
    \item revise the model if important residual structure remains.
\end{enumerate}

%%%%%%%%%%%%%%%%%%%%%%%%%%%%%%%%%%%%%%%%%%%%%%%%%%%%%%%%%%%%
\subsection{Common Errors}
\label{subsec:time-series-common-errors}
%%%%%%%%%%%%%%%%%%%%%%%%%%%%%%%%%%%%%%%%%%%%%%%%%%%%%%%%%%%%

\subsubsection{Treating Time-Series Observations as Independent}

Serial dependence changes estimator uncertainty.

Ordinary independent-sample formulas may substantially underestimate or
overestimate standard errors.

%%%%%%%%%%%%%%%%%%%%%%%%%%%%%%%%%%%%%%%%%%%%%%%%%%%%%%%%%%%%
\subsubsection{Ignoring Time Order}

Randomly rearranging a time series destroys dependence and temporal
structure.

%%%%%%%%%%%%%%%%%%%%%%%%%%%%%%%%%%%%%%%%%%%%%%%%%%%%%%%%%%%%
\subsubsection{Randomly Splitting Time-Series Training and Test Data}

Random splitting can allow future observations to influence models used
to predict the past.

Use time-respecting validation.

%%%%%%%%%%%%%%%%%%%%%%%%%%%%%%%%%%%%%%%%%%%%%%%%%%%%%%%%%%%%
\subsubsection{Confusing Trend with Autocorrelation}

A strong deterministic trend can produce large sample autocorrelations
even if residuals around the trend are independent.

Trend should be handled before interpreting stationary dependence.

%%%%%%%%%%%%%%%%%%%%%%%%%%%%%%%%%%%%%%%%%%%%%%%%%%%%%%%%%%%%
\subsubsection{Differencing Every Series Automatically}

Differencing can remove useful long-run information and can introduce
unnecessary moving-average structure.

Use it only when justified.

%%%%%%%%%%%%%%%%%%%%%%%%%%%%%%%%%%%%%%%%%%%%%%%%%%%%%%%%%%%%
\subsubsection{Failing to Difference a Unit-Root Process}

Fitting stationary ARMA models directly to strongly unit-root
nonstationary data can produce misleading inference.

%%%%%%%%%%%%%%%%%%%%%%%%%%%%%%%%%%%%%%%%%%%%%%%%%%%%%%%%%%%%
\subsubsection{Confusing Deterministic and Stochastic Trends}

A deterministic trend may be removed by regression on time.

A stochastic unit-root trend generally requires different treatment.

%%%%%%%%%%%%%%%%%%%%%%%%%%%%%%%%%%%%%%%%%%%%%%%%%%%%%%%%%%%%
\subsubsection{Using Ordinary \(t\)-Critical Values for Unit-Root Tests}

Unit-root null distributions are nonstandard.

Specialized critical values are required.

%%%%%%%%%%%%%%%%%%%%%%%%%%%%%%%%%%%%%%%%%%%%%%%%%%%%%%%%%%%%
\subsubsection{Interpreting ACF Cutoffs Too Literally}

Sample ACF and PACF values are noisy.

Finite-sample plots rarely display perfect theoretical patterns.

%%%%%%%%%%%%%%%%%%%%%%%%%%%%%%%%%%%%%%%%%%%%%%%%%%%%%%%%%%%%
\subsubsection{Ignoring Seasonality}

Seasonal dependence can remain even after ordinary differencing.

Inspect seasonal lags.

%%%%%%%%%%%%%%%%%%%%%%%%%%%%%%%%%%%%%%%%%%%%%%%%%%%%%%%%%%%%
\subsubsection{Ignoring Residual Autocorrelation}

A fitted model with strongly autocorrelated residuals has not captured
all predictable temporal structure.

%%%%%%%%%%%%%%%%%%%%%%%%%%%%%%%%%%%%%%%%%%%%%%%%%%%%%%%%%%%%
\subsubsection{Using In-Sample Fit as Forecast Accuracy}

A model can fit historical data very well and forecast poorly.

Evaluate on future or rolling-origin observations.

%%%%%%%%%%%%%%%%%%%%%%%%%%%%%%%%%%%%%%%%%%%%%%%%%%%%%%%%%%%%
\subsubsection{Assuming Narrow Parameter Intervals Mean Narrow Forecast Intervals}

Forecast uncertainty includes future random shocks in addition to
parameter uncertainty.

%%%%%%%%%%%%%%%%%%%%%%%%%%%%%%%%%%%%%%%%%%%%%%%%%%%%%%%%%%%%
\subsubsection{Ignoring Structural Breaks}

A model estimated from old observations may be irrelevant after a regime
change.

%%%%%%%%%%%%%%%%%%%%%%%%%%%%%%%%%%%%%%%%%%%%%%%%%%%%%%%%%%%%
\subsubsection{Interpreting Granger Causality as Scientific Causality}

Granger causality is based on predictive improvement from lagged
variables.

It does not by itself establish intervention-based causation.

%%%%%%%%%%%%%%%%%%%%%%%%%%%%%%%%%%%%%%%%%%%%%%%%%%%%%%%%%%%%
\subsubsection{Regressing Unrelated Nonstationary Series}

Spurious regression can create apparently strong relationships between
unrelated trending processes.

%%%%%%%%%%%%%%%%%%%%%%%%%%%%%%%%%%%%%%%%%%%%%%%%%%%%%%%%%%%%
\subsubsection{Ignoring Parameter Uncertainty in Forecasts}

Simple theoretical forecast intervals often condition on estimated
parameters as if known.

Finite-sample forecast uncertainty may therefore be larger.

%%%%%%%%%%%%%%%%%%%%%%%%%%%%%%%%%%%%%%%%%%%%%%%%%%%%%%%%%%%%
\subsection{Applications}
\label{subsec:time-series-applications}
%%%%%%%%%%%%%%%%%%%%%%%%%%%%%%%%%%%%%%%%%%%%%%%%%%%%%%%%%%%%

\begin{applicationbox}

\textbf{Environmental Science.}
Temperature, rainfall, dissolved oxygen, salinity, and sensor
measurements are often recorded repeatedly through time and exhibit
seasonal and serial dependence.

\medskip

\textbf{Economics.}
GDP, inflation, employment, interest rates, and other economic indicators
are modeled using ARIMA, VAR, cointegration, and state-space methods.

\medskip

\textbf{Finance.}
Asset prices, returns, and volatility require models for dependence,
nonstationarity, and time-varying variance.

\medskip

\textbf{Engineering.}
Sensor streams, vibration measurements, control systems, and reliability
signals are naturally modeled as stochastic processes.

\medskip

\textbf{Business.}
Sales, demand, inventory, call volume, and web traffic frequently require
trend and seasonal forecasting.

\medskip

\textbf{Medicine.}
Repeated physiological measurements, monitoring signals, and disease
trajectories often exhibit temporal dependence.

\medskip

\textbf{Epidemiology.}
Disease counts and incidence rates through time can contain trends,
seasonality, interventions, and structural changes.

\medskip

\textbf{Operations.}
Forecasts of arrivals, demand, workload, and resource usage support
scheduling and inventory decisions.

\end{applicationbox}

%%%%%%%%%%%%%%%%%%%%%%%%%%%%%%%%%%%%%%%%%%%%%%%%%%%%%%%%%%%%
\subsection{Exercises}
\label{subsec:time-series-exercises}
%%%%%%%%%%%%%%%%%%%%%%%%%%%%%%%%%%%%%%%%%%%%%%%%%%%%%%%%%%%%

\begin{exercise}[1]
Define a time series.
\end{exercise}

\begin{exercise}[1]
Define lag.
\end{exercise}

\begin{exercise}[1]
Define weak stationarity.
\end{exercise}

\begin{exercise}[1]
Define the autocovariance function.
\end{exercise}

\begin{exercise}[1]
Define the autocorrelation function.
\end{exercise}

\begin{exercise}[1]
Define white noise.
\end{exercise}

\begin{exercise}[1]
Define an AR(1) process.
\end{exercise}

\begin{exercise}[1]
Define an MA(1) process.
\end{exercise}

\begin{exercise}[1]
Define an ARMA(\(p,q\)) process.
\end{exercise}

\begin{exercise}[1]
Explain the meaning of ARIMA(\(p,d,q\)).
\end{exercise}

\begin{exercise}[1]
Define first differencing.
\end{exercise}

\begin{exercise}[1]
Define seasonal differencing.
\end{exercise}

\begin{exercise}[1]
Define the PACF conceptually.
\end{exercise}

\begin{exercise}[1]
Define a forecast error.
\end{exercise}

\begin{exercise}[2]
For

\[
X_t
=
0.4X_{t-1}
+
\varepsilon_t,
\]

determine whether the AR(1) process is stationary.
\end{exercise}

\begin{exercise}[2]
For

\[
X_t
=
1.2X_{t-1}
+
\varepsilon_t,
\]

determine whether the AR(1) process is stationary.
\end{exercise}

\begin{exercise}[2]
Suppose

\[
X_t
=
2
+
0.5X_{t-1}
+
\varepsilon_t.
\]

Find the stationary mean.
\end{exercise}

\begin{exercise}[2]
Suppose an AR(1) has

\[
\phi=0.8
\]

and

\[
\sigma_\varepsilon^2=9.
\]

Find the stationary variance.
\end{exercise}

\begin{exercise}[2]
For the previous AR(1), compute

\[
\rho(1),
\quad
\rho(2),
\quad
\rho(3).
\]
\end{exercise}

\begin{exercise}[2]
Suppose

\[
X_t
=
\varepsilon_t
+
0.5\varepsilon_{t-1}.
\]

Find the variance.
\end{exercise}

\begin{exercise}[2]
For the previous MA(1), compute \(\rho(1)\).
\end{exercise}

\begin{exercise}[2]
For the previous MA(1), find \(\rho(2)\).
\end{exercise}

\begin{exercise}[2]
Classify a random walk as an ARIMA model.
\end{exercise}

\begin{exercise}[2]
Write

\[
X_t-X_{t-12}
\]

using the backshift operator.
\end{exercise}

\begin{exercise}[2]
Expand

\[
(1-B)^2X_t.
\]
\end{exercise}

\begin{exercise}[2]
For a monthly series with annual seasonality, state the seasonal period.
\end{exercise}

\begin{exercise}[2]
If the AR(1) stationary mean is \(\mu=20\), \(\phi=0.7\), and
\(X_t=25\), find the one-step forecast.
\end{exercise}

\begin{exercise}[2]
For the previous model, find the two-step forecast.
\end{exercise}

\begin{exercise}[2]
For a random walk with innovation variance \(4\), find the variance of a
five-step forecast error.
\end{exercise}

\begin{exercise}[2]
Suppose simple exponential smoothing uses

\[
\alpha=0.3,
\]

with

\[
X_t=20,
\qquad
\ell_{t-1}=18.
\]

Find \(\ell_t\).
\end{exercise}

\begin{exercise}[3]
Derive the stationary mean of an AR(1) process.
\end{exercise}

\begin{exercise}[3]
Derive the stationary variance of an AR(1) process.
\end{exercise}

\begin{exercise}[3]
Derive

\[
\rho(h)
=
\phi^{|h|}
\]

for a stationary AR(1).
\end{exercise}

\begin{exercise}[3]
Derive the infinite moving-average representation of a stationary AR(1).
\end{exercise}

\begin{exercise}[3]
Derive the variance of an MA(1) process.
\end{exercise}

\begin{exercise}[3]
Derive the lag-one autocorrelation of an MA(1).
\end{exercise}

\begin{exercise}[3]
Show that an MA(1) has zero autocorrelation beyond lag \(1\).
\end{exercise}

\begin{exercise}[3]
Explain the difference between stationarity and invertibility.
\end{exercise}

\begin{exercise}[3]
Derive the Yule--Walker equations for an AR(2) process.
\end{exercise}

\begin{exercise}[3]
Explain why the PACF of an AR(\(p\)) process cuts off after lag \(p\).
\end{exercise}

\begin{exercise}[3]
Explain why the ACF of an MA(\(q\)) process cuts off after lag \(q\).
\end{exercise}

\begin{exercise}[3]
Explain why differencing removes a random-walk stochastic trend.
\end{exercise}

\begin{exercise}[3]
Compare deterministic detrending and stochastic differencing.
\end{exercise}

\begin{exercise}[3]
Explain why a unit-root test does not use the usual \(t\)-distribution.
\end{exercise}

\begin{exercise}[3]
Explain why two unrelated random walks can produce spurious regression.
\end{exercise}

\begin{exercise}[3]
Derive the one-step AR(1) forecast.
\end{exercise}

\begin{exercise}[3]
Derive the \(h\)-step AR(1) forecast.
\end{exercise}

\begin{exercise}[3]
Derive the AR(1) \(h\)-step forecast error variance.
\end{exercise}

\begin{exercise}[3]
Explain why random-walk forecast uncertainty grows with horizon.
\end{exercise}

\begin{exercise}[3]
Explain why time-series validation should preserve temporal ordering.
\end{exercise}

\begin{exercise}[3]
Compare expanding-window and rolling-window forecasting.
\end{exercise}

\begin{exercise}[3]
Show that exponential-smoothing weights decay geometrically.
\end{exercise}

\begin{exercise}[3]
Explain when additive seasonality is more suitable than multiplicative
seasonality.
\end{exercise}

\begin{exercise}[3]
Explain why residual autocorrelation indicates model inadequacy.
\end{exercise}

\begin{exercise}[3]
Explain the difference between an intervention variable and ordinary
random noise.
\end{exercise}

\begin{exercise}[3]
Explain the distinction between VAR and univariate AR models.
\end{exercise}

\begin{exercise}[3]
Explain the meaning of cointegration.
\end{exercise}

\begin{exercise}[3]
Explain why Granger causality is not equivalent to scientific causality.
\end{exercise}

\begin{exercise}[3]
Show that white noise has a constant spectral density.
\end{exercise}

\begin{exercise}[4]
Derive the stationary solution of a general AR(\(p\)) process under its
root condition.
\end{exercise}

\begin{exercise}[4]
Study the causal MA(\(\infty\)) representation of ARMA models.
\end{exercise}

\begin{exercise}[4]
Derive invertibility conditions for MA(\(q\)) processes.
\end{exercise}

\begin{exercise}[4]
Derive the Yule--Walker estimator for AR(\(p\)) parameters.
\end{exercise}

\begin{exercise}[4]
Study maximum likelihood estimation for Gaussian ARMA models.
\end{exercise}

\begin{exercise}[4]
Investigate innovations algorithms for ARMA models.
\end{exercise}

\begin{exercise}[4]
Study Dickey--Fuller and augmented Dickey--Fuller unit-root tests.
\end{exercise}

\begin{exercise}[4]
Investigate KPSS-type stationarity tests.
\end{exercise}

\begin{exercise}[4]
Study seasonal unit roots.
\end{exercise}

\begin{exercise}[4]
Derive properties of seasonal ARIMA models.
\end{exercise}

\begin{exercise}[4]
Study forecast distributions with parameter uncertainty.
\end{exercise}

\begin{exercise}[4]
Investigate optimal forecast combinations.
\end{exercise}

\begin{exercise}[4]
Study the Kalman filter equations for a linear Gaussian state-space
model.
\end{exercise}

\begin{exercise}[4]
Derive the local-level model as a state-space process.
\end{exercise}

\begin{exercise}[4]
Study vector autoregression estimation.
\end{exercise}

\begin{exercise}[4]
Investigate impulse-response functions in VAR models.
\end{exercise}

\begin{exercise}[4]
Study Granger-causality tests in VAR systems.
\end{exercise}

\begin{exercise}[4]
Study cointegration and error-correction models.
\end{exercise}

\begin{exercise}[4]
Investigate the Engle--Granger cointegration procedure.
\end{exercise}

\begin{exercise}[4]
Study the Johansen cointegration framework.
\end{exercise}

\begin{exercise}[4]
Derive the Wiener--Khinchin relationship between autocovariance and
spectral density.
\end{exercise}

\begin{exercise}[4]
Study the periodogram and its statistical properties.
\end{exercise}

\begin{exercise}[4]
Investigate ARCH and GARCH models.
\end{exercise}

\begin{exercise}[4]
Study change-point and structural-break models.
\end{exercise}

%%%%%%%%%%%%%%%%%%%%%%%%%%%%%%%%%%%%%%%%%%%%%%%%%%%%%%%%%%%%
\subsection{Conceptual Summary}
\label{subsec:time-series-summary}
%%%%%%%%%%%%%%%%%%%%%%%%%%%%%%%%%%%%%%%%%%%%%%%%%%%%%%%%%%%%

A time series is an ordered stochastic process

\[
\boxed{
\{X_t\}.
}
\]

For a weakly stationary process,

\[
\boxed{
\mathbb{E}[X_t]=\mu,
}
\]

and

\[
\boxed{
\operatorname{Cov}
(
X_t,
X_{t-h}
)
=
\gamma(h).
}
\]

The autocorrelation function is

\[
\boxed{
\rho(h)
=
\frac{
\gamma(h)
}{
\gamma(0)
}.
}
\]

White noise satisfies

\[
\boxed{
\gamma(h)=0
\qquad
h\neq0.
}
\]

A stationary AR(1) process is

\[
\boxed{
X_t
=
c
+
\phi X_{t-1}
+
\varepsilon_t,
\qquad
|\phi|<1.
}
\]

Its mean is

\[
\boxed{
\mu
=
\frac{
c
}{
1-\phi
},
}
\]

and its autocorrelation is

\[
\boxed{
\rho(h)
=
\phi^{|h|}.
}
\]

An MA(1) process is

\[
\boxed{
X_t
=
\varepsilon_t
+
\theta\varepsilon_{t-1},
}
\]

with autocorrelation

\[
\boxed{
\rho(1)
=
\frac{
\theta
}{
1+\theta^2
}
}
\]

and

\[
\boxed{
\rho(h)=0
\qquad
|h|>1.
}
\]

ARIMA models combine autoregression, differencing, and moving-average
structure:

\[
\boxed{
ARIMA(p,d,q).
}
\]

Forecasting uses conditional expectation:

\[
\boxed{
\widehat{X}_{t+h\mid t}
=
\mathbb{E}
[
X_{t+h}
\mid
\mathcal{F}_t
].
}
\]

Time series analysis therefore combines

\[
\boxed{
\text{probability}
+
\text{stochastic processes}
+
\text{regression}
+
\text{difference equations}
+
\text{forecasting}.
}
\]

%%%%%%%%%%%%%%%%%%%%%%%%%%%%%%%%%%%%%%%%%%%%%%%%%%%%%%%%%%%%
\subsection{Section Connections}
\label{subsec:time-series-connections}
%%%%%%%%%%%%%%%%%%%%%%%%%%%%%%%%%%%%%%%%%%%%%%%%%%%%%%%%%%%%

\chapterconnection{
Time Series Analysis applies the stochastic-process ideas developed in
Chapter~\ref{chap:probability} to ordered statistical observations.
Autoregressive models are stochastic difference equations, while ARIMA
models combine autoregressive dynamics, differencing, and innovations.
Regression models can be extended to include serially correlated errors,
lagged predictors, and interventions. State-space models connect time
series with dynamical systems and filtering, while vector autoregressions
extend the framework to multivariate processes. The next section,
Spatial Statistics, studies dependence indexed by geographic or spatial
location rather than primarily by time.
}

%%%%%%%%%%%%%%%%%%%%%%%%%%%%%%%%%%%%%%%%%%%%%%%%%%%%%%%%%%%%
% SECTION SOURCE TRACKING
%%%%%%%%%%%%%%%%%%%%%%%%%%%%%%%%%%%%%%%%%%%%%%%%%%%%%%%%%%%%

%------------------------------------------------------------
% 8.15 Time Series Analysis
%------------------------------------------------------------
%
% Source basis:
%   - Standard time-series theory.
%   - Standard stochastic-process and forecasting methods.
%   - Standard ARIMA and state-space concepts.
%
% Used for:
%   - time-series structure;
%   - trend, seasonality, cycles, and irregular variation;
%   - additive and multiplicative decomposition;
%   - lag;
%   - weak and strict stationarity;
%   - autocovariance;
%   - autocorrelation;
%   - sample ACF;
%   - white noise;
%   - random walks;
%   - differencing;
%   - seasonal differencing;
%   - backshift operators;
%   - autoregressive models;
%   - AR(1) mean, variance, and ACF;
%   - AR stationarity conditions;
%   - Yule--Walker equations;
%   - moving-average models;
%   - MA(1) variance and ACF;
%   - invertibility;
%   - ARMA models;
%   - PACF;
%   - ARIMA models;
%   - seasonal ARIMA;
%   - deterministic and stochastic trends;
%   - unit roots;
%   - spurious regression;
%   - transformations;
%   - forecasting;
%   - forecast errors and intervals;
%   - exponential smoothing;
%   - Holt and Holt--Winters methods;
%   - time-series regression;
%   - dynamic regression;
%   - intervention analysis;
%   - transfer functions;
%   - residual diagnostics;
%   - Ljung--Box tests;
%   - AIC and BIC;
%   - state-space models;
%   - Kalman filter preview;
%   - multivariate time series;
%   - VAR models;
%   - cross-covariance;
%   - Granger causality preview;
%   - cointegration;
%   - error-correction models;
%   - frequency-domain analysis;
%   - spectral density;
%   - periodograms;
%   - structural breaks;
%   - ARCH and GARCH previews;
%   - applications and exercises.
%
% Notes:
%   - Spatial Statistics is developed next in Section 8.16.
%   - Stochastic Processes in Chapter 7 provides the underlying
%     probability theory.
%   - Differential Equations and Dynamical Systems later develops
%     deterministic and stochastic dynamical models more deeply.
%   - Computational Statistics later develops numerical forecasting,
%     simulation, and state-estimation methods.
%
%%%%%%%%%%%%%%%%%%%%%%%%%%%%%%%%%%%%%%%%%%%%%%%%%%%%%%%%%%%%